\documentclass[11pt,reqno]{amsart}

\usepackage[dvipsnames]{xcolor}
\usepackage{amsmath}
\usepackage{amsthm}
\usepackage{graphicx}
\usepackage{upgreek}
\usepackage{hyperref}
\usepackage{mathrsfs}
\usepackage{bm,bbm}
\usepackage{amsfonts}
\usepackage{amssymb}
\usepackage{csquotes}
\usepackage{stmaryrd}
\usepackage{ulem}
\usepackage{tikz-cd}
\usepackage{cancel}
\usepackage{siunitx}
\usepackage{algorithm}
\usepackage{algpseudocode}
\usepackage{caption}
\usepackage{subcaption}
\usepackage{booktabs}
\usepackage{enumerate}
\usepackage{soul}
\usepackage[margin=1.1in]{geometry}
\parindent=.25in


\numberwithin{equation}{section}

\theoremstyle{plain}
\newtheorem{theorem}{Theorem}[section]
\newtheorem{definition}[theorem]{Definition}
\newtheorem{proposition}[theorem]{Proposition}
\newtheorem{lemma}[theorem]{Lemma}
\newtheorem{question}[theorem]{Question}
\newtheorem{remark}[theorem]{Remark}
\newtheorem{example}[theorem]{Example}
\newtheorem{corollary}[theorem]{Corollary}

\newcommand{\thistheoremname}{}
\newtheorem{genericthm}[theorem]{\thistheoremname}
\newenvironment{assumption}[1]
  {\renewcommand{\thistheoremname}{#1}%
   \begin{genericthm}}
  {\end{genericthm}}

\newtheorem*{genericthm*}{\thistheoremname}
\newenvironment{assumption*}[1]
  {\renewcommand{\thistheoremname}{#1}
   \begin{genericthm*}}
  {\end{genericthm*}}

% Character macros
\newcommand{\Z}{\mathbb{Z}}
\newcommand{\R}{\mathbb{R}}
\newcommand{\reals}{\mathbb{R}}
\newcommand{\bR}{\mathbf{R}}
\newcommand{\bP}{\mathbf{P}}
\newcommand{\bA}{\mathbb{A}}
\newcommand{\bQ}{\mathbf{Q}}
\newcommand{\Q}{\mathbb{Q}}
\newcommand{\C}{\mathbb{C}}
\newcommand{\X}{\mathbb{X}}
\newcommand{\T}{\mathbb{T}}
\newcommand{\F}{\mathbb{F}}
\newcommand{\sF}{\mathcal{F}}
\renewcommand{\L}{\mathcal{L}}
\renewcommand{\P}{\mathbb{P}}
\newcommand{\sP}{\mathcal{P}}
\newcommand{\rcp}{\mathbf{\widehat{P}}}
\newcommand{\sA}{\mathcal{A}}
\newcommand{\sG}{\mathcal{G}}
\newcommand{\sS}{\mathcal{S}}
\newcommand{\sI}{\mathcal{I}}
\newcommand{\sX}{\mathcal{X}}
\newcommand{\sU}{\mathcal{U}}
\newcommand{\sZ}{\mathcal{Z}}
\newcommand{\sC}{\mathcal{C}}
\newcommand{\sT}{\mathcal{T}}
\newcommand{\sB}{\mathcal{B}}
\newcommand{\sN}{\mathcal{N}}
\newcommand{\bX}{\mathbf{X}}
\newcommand{\sH}{\mathcal{H}}
\newcommand{\sM}{\mathcal{M}}
\newcommand{\sL}{\mathcal{L}}
\newcommand{\N}{\mathbb{N}}
\newcommand{\h}{\mathbb{H}}
\renewcommand{\S}{\mathbb{S}}
\newcommand{\D}{\mathbb{D}}
\newcommand{\E}{\mathbb{E}}

% Other miscellaneous macros
\newcommand{\lbrs}[1]{\llbracket #1 \rrbracket}
\newcommand{\io}{\text{ i.o.}}
\newcommand{\ind}{\mathbbm{1}}
\newcommand{\sign}{\text{sign}}
\newcommand{\Var}{\text{Var}}
\newcommand{\Span}{\text{span}}
\newcommand{\Img}{\text{Im }}
\newcommand{\Rl}{\text{Re}}
\newcommand{\Cov}{\text{Cov}}
\newcommand{\intr}{\text{Int}}
\newcommand{\Syl}{\text{Syl}}
\newcommand{\opnorm}{_{\text{op}}}
\newcommand{\trace}{{\text{tr}}}
\newcommand{\mc}[1]{\mathcal{#1}}
\newcommand{\comp}[1]{\langle #1 \rangle}
\newcommand{\bmid}{\;\bigg| \;}
\newcommand{\argmin}{\mathop{\rm argmin}}
\newcommand{\argmax}{\mathop{\rm argmax}}


\newcommand{\AC}[1]{\texttt{{{\textcolor{red}{(AC:\ #1)}}}}}
\newcommand{\JS}[1]{\texttt{{{\textcolor{magenta}{(JS:\ #1)}}}}}
\newcommand{\JH}[1]{\texttt{{{\textcolor{blue}{(JH:\ #1)}}}}}
\newcommand{\credb}[2]{\cancel{\textcolor{red}{{#1}}} \textcolor{blue}{#2}}


\title{Convergence of the Deep Galerkin Method for Mean Field Control Problems}


\author{William Hofgard}
\address{Department of Mathematics, Stanford University, Stanford}
\email{whofgard@stanford.edu}

\author{Jingruo Sun}
\address{Department of Management Science and Engineering, Stanford University, Stanford}
\email{jingruo@stanford.edu}


\author{Asaf Cohen}
\address{Department of Mathematics, University of Michigan, Ann Arbor}
\email{shloshim@gmail.com}
\urladdr{https://sites.google.com/site/asafcohentau/}

\thanks{The research of A. Cohen and W. Hofgard was supported by the National Science Foundation (DMS 2006305). The research of J. Sun was supported by the REU funds of the Department of Mathematics at the University of Michigan.}
\def\namedlabel#1#2{\begingroup
    #2%
    \def\@currentlabel{#2}%
    \phantomsection\label{#1}\endgroup
}

\date{\today}

\begin{document}
\maketitle
\begin{abstract}
We establish the convergence of the deep Galerkin method (DGM), a deep learning-based scheme for solving high-dimensional nonlinear PDEs, for Hamilton-Jacobi-Bellman (HJB) equations that arise from the study of mean field control problems (MFCPs). Based on a recent characterization of the value function of the MFCP as the unique viscosity solution of an HJB equation on the simplex, we establish both an existence and convergence result for the DGM. First, we show that the loss functional of the DGM can be made arbitrarily small given that the value function of the MFCP possesses sufficient regularity. Then, we show that if the loss functional of the DGM converges to zero, the corresponding neural network approximators must converge uniformly to the true value function on the simplex. We also provide numerical experiments demonstrating the DGM's ability to generalize to high-dimensional HJB equations.
\end{abstract}
\noindent{\bf Keywords:} mean field control problems, deep Galerkin method, neural networks, PDEs 

\noindent{\bf AMS Classification:} 
91A07,  	%Games with infinitely many players
35Q89,  	%PDEs in connection with mean field game theory
68T07,  	%Artificial neural networks and deep learning
49L12,  	%Hamilton-Jacobi equations in optimal control and differential games
49N10, 	%Linear-quadratic optimal control problems
35A35,  	%Theoretical approximation in context of PDEs
60J27 %Continuous-time Markov processes on discrete state spaces

\tableofcontents

\section{Introduction}
We consider the convergence of a deep learning-based numerical method for solving Hamilton--Jacobi--Bellman (HJB) equations that arise from the study of mean field control problems (MFCPs). The specific class of HJB equations that we consider are first-order, nonlinear PDEs on the $(d-1)$--dimensional simplex $S_d$, of the following form:
\begin{align}
\label{eq:mfcp-hjb}
\begin{split}
    &-\partial_t V(t, m) + \sum_{i \in \lbrs{d}} m_i H^i(t, m, D^i V(t,m)) = 0,\\
    &V(T, m) = \sum_{i \in \lbrs{d}} m_i g^i(m).
\end{split}
\end{align}
Above, $\lbrs{d} := \{1, \ldots, d\}$, $V : [0,T] \times S_d \to \R$ is the {\it value function} (optimal cost) of an associated MFCP, $H^i$ is a Hamiltonian (to be introduced below in the context of a MFCPs), and $D^i$ denotes a derivative along the simplex, namely  $D^i V(t,m)=(\partial_{m_j - m_i} V(t,m))_{j \in \lbrs{d}}$. Note that on the simplex, directional derivatives are only permitted in the directions $e_j - e_i$, where $e_i$ denotes the $i$th standard basis vector in $\R^d$. It is Equation~\eqref{eq:mfcp-hjb} that we focus on solving numerically due to the connection between the MFCP and the $N$-agent optimization problem. 

Our numerical approach is primarily based on the deep Galerkin method (DGM), as first introduced by \cite{dgm}. In particular, we solve the above class of PDEs using a deep learning-based algorithm motivated by the classical finite element method. However, instead of finding a basis of functions to approximate solutions as in the classical method, we utilize a sufficiently rich class of neural networks to approximate solutions to Equation~\eqref{eq:mfcp-hjb}. In addition to proving the convergence of our numerical scheme under the assumptions outlined in Section~\ref{sec:hjb-mfcp-setup}, we provide numerical experiments that demonstrate the validity of this approach. First, however, we introduce the MFCP model in Section~\ref{subsec:mfcp}, describe the current field of solving high-dimensional PDE using deep learning in Section~\ref{subsec:num-sol-hpde}, and outline our main results in Section~\ref{subsec:main-results}.

\subsection{Mean Field Control Problems}
\label{subsec:mfcp}
\subsubsection{Literature review}
MFCPs are limiting models for cooperative games with a finite, but large, number of interacting agents attempting to minimize a common cost. Namely, the study of MFCPs is motivated by situations in which a large number of agents aim to achieve a common goal. For instance, applications arise in the control of autonomous vehicles and drones \cite{drones7070476, mot-pha2022CONTROL}, efficient real-time streaming between devices \cite{7218597-2015}, decentralized and centralized crowd motion \cite{crowd-motion-1, crowd-motion-2}, analysis of wireless local-area networks (WLANs) \cite{duffy2010meanCONTROL, bor-sun2012CONTROL}, among many other related applications. The inception of limiting models for many-agent games can be traced back to Aumann's seminal work in 1964 \cite{aumann1964markets} and Schmeidler's influential contribution in 1973 \cite{schmeidler1973equilibrium}.  These works dealt with one-shot games. In the context of stochastic dynamical games, mean field game (MFG) theory was initially developed by Lasry and Lions \cite{lasry-lions-mfg} and independently by Huang, Caines, and Malhamé \cite{huang-caines-malhame-mfg}. In our discourse, we distinguish between MFCPs and MFGs, which are frameworks for approximating, respectively, cooperative and non-cooperative games involving numerous players. MFCPs represent optimal control problems for McKean--Vlasov dynamics, where the value function (optimal cost) is often delineated by a HJB equation. Conversely, in MFGs, the Nash equilibrium is approximated through a fixed point of a best response map, and the value function (i.e., the cost under the proper equilibrium concept in the MFG) is characterized by what is called the master equation. 

Recent work by Pham and Wei derived a dynamic programming principle and HJB equation for the MFCP by framing it as a deterministic control problem of a Fokker--Planck equation; see \cite{Pham2017, Pham2018}. Even more recently, work by Cecchin \cite{mfcp} establishes important existence and uniqueness results for the HJB equation corresponding to the finite-state MFCP we consider here. For a comprehensive overview of MFG theory, MFCPs, and their relationship to stochastic games and control problems, see \cite{Bensoussan2013, Carmona2018_I, Carmona2018_II}.  


Given current advances in deep learning, there has been recent progress towards solving both MFGs and MFCPs numerically using deep learning \cite{ml-mfcp, deep-learning-mfcp}. Since 2020, several review articles have again brought up the potential for solving MFGs and MFCPs with deep learning methods, including the DGM and deep BSDE, but they note that so far that no convergence guarantee has been provided for these methods; see \cite{mfc-mfcp-review, scalable-rl-mfg}. Similarly, \cite{actor-critic-bsde} provides a broad sampling of applications of deep BSDE-like algorithms to common problems in stochastic control, again without providing a convergence guarantee for their proposed method. Building on the theoretical results from \cite{mfcp}, we aim to provide a novel pathway towards solving MFCPs using deep learning, complete with a robust convergence guarantee.





\subsubsection{MFCP model} The $N$-agent optimization problem, which limits to the MFCP as $N$ approaches infinity, is a continuous-time, finite state optimization problem with $N$ agents whose state-process is $(\mathbf{X}_s = (X^1_s, \ldots, X^N_s))_{s\in[t_0,T]}$, where each coordinate takes values in the state space $\lbrs{d} = \{1, \ldots, d\}$; $[t_0,T]$ is the deterministic time interval of the problem. The agents' state process $(\mathbf{X}_t)_{t\in[t_0,T]}$ follow the dynamics given by the continuous-time Markov chain, where each of its coordinates follows 
\begin{align}
\label{eq:n-agent}
    \P(X_{s+h}^k = j \mid \mathbf{X}_s = \mathbf{x}) = \beta^k(s, \mathbf{x};j)h + o(h)\qquad\text{as \qquad $h\to0^+$},\qquad \forall j\in\lbrs{d},\; j\ne x^k.
\end{align}
The agents choose feedback controls $(\beta^1, \ldots, \beta^N)$ in order to minimize the common cost
\begin{align}
    J^N(t_0, m^N_{t_0}, \beta^1, \ldots, \beta^N) := \frac{1}{N} \sum_{k=1}^N\E\left[\int_{t_0}^T f(s, X_s^k, \beta^k(s, \mathbf{X}_s;\cdot), \mu_s^N) ds + g(X^k_T, \mu_T^N)\right] \label{eq:n-agent-cost},
\end{align}
where $\mu^N_s$ is the empirical distribution of the $N$ agents at time $s$, satisfying the predetermined initial condition $\mu^N_{t_0}=m^N_{t_0}$, $f$ denotes the running cost of the stochastic control problem while $g$ denotes the terminal cost. The value of $f(s,i,a,m)$ is independent of $a_i$ for each $i \in \lbrs{d}$ because there is no transition rate between any state $i \in \lbrs{d}$ and itself. 

In practice, finding the optimal control in $N$-agent stochastic control problems quickly becomes intractable as $N$ grows. As a result, the MFCP is of theoretical and practical interest in stochastic control. Before presenting it we provide a short motivation for its structure. More precisely, as $N$ approaches infinity, the $N$-agent optimization problem resembles an optimization problem involving a single {\it representative} agent whose dynamics follow a controlled continuous-time Markov chain of the form:
\begin{align}
\label{eq:mfcp-process}
    \P(X_{s+h} = j \mid X_s = i) = \alpha_{i,j}(s) h + o(h)\qquad\text{as}\qquad h\to0^+,
\end{align}
where the feedback control of the representative agent $\alpha_{i,j}(s)$ is a measurable, Markovian function, standing for the transition rate of the above process from state $i$ to $j$, both in $\lbrs{d}$.
Throughout this paper, we consider only rate matrices satisfying $\alpha_{i,j}(s) \in [0, M]$ for some constant $M > 0$ and all $i\ne j, i,j \in\lbrs{d}$, $s\in[t_0,T]$, and we refer to them as {\it admissible controls}. We denote the set of admissible controls by $\mathcal{A}_{t_0}$. Also, $\alpha_i$ refers to the $i$th row of the transition matrix $\alpha$. 
At this point, the reader may wonder about the gap between the richness of the feedback controls allowed in the $N$-agent problem versus the MFCP. We posit that this apparent disparity is unfounded. Given the symmetry among agents, a logical step involves substituting their transition rates with functions of the form $\bar \beta^k(s,X^k_s,\mu^N_s;\cdot)$. Furthermore, guided by the propagation of chaos result, we anticipate the convergence of $(\mu^N_t)_{t\in [t_0,T]}$ toward a deterministic flow of measures, approximating the distributions of each coordinate in $\mathbf{X}$. This deterministic framework implies that, in the limit, we can anticipate the control's dependence on both the current time $t$ and the present state $i$, aligning with our prior considerations.

Similar to the $N$-agent case, the representative agent aims to minimize the cost
\begin{align}
    J(t_0, m_{t_0}, \alpha) := \E\left[\int_{t_0}^T f(s, X_s, \alpha_{X_s}(s), \text{Law}(X_s))ds + g(X_T, \text{Law}(X_T))\right] \label{eq:mfcp-cost},
\end{align}
where the initial distribution of the process $X$: $\text{Law}(X_{t_0}) = m_{t_0}\in S_d$ is predetermined. This problem can be alternatively described as a deterministic control problem in terms of the Fokker--Planck equation for $(\mu_s := \text{Law}(X_s))_{s\in[t_0,T]}$. In particular, for a given admissible control $\alpha$, $\mu$ uniquely solves the ordinary differential equation (ODE):
\begin{align}
\label{eq:forward-eq}
\begin{split}
    &\frac{d}{dt}\mu_t^i = \sum_{j \in \lbrs{d}} \left(\mu_t^j \alpha_{j,i}(s) - \mu_t^i \alpha_{i,j}(s)\right),\qquad i\in\lbrs{d}, \quad t\in[t_0,T], \\
    &\mu_{t_0} = m_{t_0},
\end{split}
\end{align}
and the cost functional for the deterministic control problem becomes
\begin{align}
    \label{eq:det-mfcp-cost}
    J(t_0, m_{t_0},\alpha) := \int_{t_0}^{T} \sum_{i \in \lbrs{d}} f(s, i, \alpha_{i}(s), \mu_s) \mu_s^i ds + \sum_{i \in \lbrs{d}} g^i(\mu_T) \mu_T^i.
\end{align}
Allowing for the initial time of the problem $t_0$ to be in $[0,T]$, the {\it value function} $V:[0,T]\times S_d\to\R$ of the deterministic control problem is then defined by
\begin{align}
    \label{eq:value-function}
    V(t, m) = \inf_{\alpha \in \mathcal{A}_{t}} J(t, \alpha, m).
\end{align}
By standard optimal control arguments via an appropriate dynamic programming principle, this gives rise to the HJB equation \eqref{eq:mfcp-hjb}, where the {\it PDE-Hamiltonian} $H : [0, T] \times S_d \times \R^d \to \R$ is defined by $H(t, m, z) = \sum_{i \in \lbrs{d}} m_i H^i(t, m, z)$, with
\begin{align}
\label{eq:hamiltonian}
    H^i(t,m,z) := \sup_{a \in [0, M], \;k\in\lbrs{d}\setminus\{i\}} \Big(-\sum_{k\in\lbrs{d}\setminus\{i\}}a_k z_k - f(t,i,a,m)\Big).
\end{align}
Note that we refer to both $H^i$ and $H$ as Hamiltonians. To avoid confusion, we occasionally employ the term PDE-Hamiltonian specifically for $H$. Under the assumptions presented in Section~\ref{subsec:assumptions} below, the above equation has a unique maximizer \cite{unique-optimal-control-mfg}, which we denote by $a^*(t, i, m, z)$ for each $i \in \lbrs{d}$. In turn, the \textit{unique} optimal control for the MFCP is given by
\begin{align}
    \label{eq:optimal-control-mfcp}
    \alpha^*_{i,j}(t,m) = a_j^*(t, i,  m, D^i V(t, m)),
\end{align}
where the directional derivative of the value function $D^i V(t, m)$ on the simplex $S^d$ is as defined in Equation~\eqref{eq:mfcp-hjb}. Additionally, the dynamics of the process $\mu$ under the optimal control satisfy the dynamics in Equation~\eqref{eq:forward-eq} upon replacing the control therein with the optimal control from Equation~\eqref{eq:optimal-control-mfcp} above \cite{mfcp}.

Importantly, recent work rigorously establishes the connection between this $N$-agent stochastic control problem and the corresponding MFCP. In particular, Cecchin \cite{mfcp} formulates the MFCP in terms of the deterministic optimal control of a Fokker--Planck equation as in Equations~\eqref{eq:forward-eq} and \eqref{eq:det-mfcp-cost} above. The resulting deterministic control problem yields a dynamic programming principle and the HJB equation \eqref{eq:mfcp-hjb} solved by the value function of the MFCP.  By formulating the $N$-agent optimization problem in a similar manner, one obtains two HJB equations: one for the $N$-agent optimization problem and one for the MFCP. Both HJB equations have unique viscosity solutions, as shown in \cite{mfcp}. In turn, \cite{mfcp} establishes an explicit rate of convergence between $V^N$, the value function for the $N$-agent problem, and $V$, the value function for the MFCP. In particular, a convergence rate of $1/\sqrt{N}$ is established. In a similar vein, Kolokoltsov \cite{kolokoltsov-conv} obtains a convergence rate of $1/N$ under additional regularity assumptions. Notably, \cite{mfcp} places no regularity assumptions on the value function $V$ (other than Lipschitz continuity), and in the most general case, does not require convexity of the running or terminal costs of the MFCP. Given the relationship between the $N$-agent problem and the MFCP, by approximately solving the HJB equation for the MFCP in this paper, we essentially provide a method for approximately solving the $N$-agent problem.

\subsection{Numerical Solutions to High-Dimensional PDE}
\label{subsec:num-sol-hpde}
We aim to construct a numerical scheme for efficiently solving the HJB equation associated with the MFCP. In particular, as the dimension $d$ increases, the so-called ``curse of dimensionality'' prevents standard numerical schemes (e.g., Monte Carlo methods, mesh-based algorithms, etc.) from solving the HJB equation in a tractable manner. However, recent advancements in deep learning present promising options for solving high-dimensional, non-linear PDE such as the HJB equation in Equation~\eqref{eq:mfcp-hjb}.

Two leading methods have been presented for parabolic PDEs that resemble the HJBs equation: the deep Galerkin method (DGM) and deep backward stochastic differential equations (BSDE). The primary focus of this paper is the DGM. This approach, first introduced in \cite{dgm}, models itself after classical finite element methods for solving low-dimensional PDEs. However, the DGM is a mesh-free method; instead of creating basis functions that approximate the solution to a PDE from a mesh, the DGM utilizes neural network approximations that only depend on the parameters, the architecture, and the activation function between layers of the network. The loss functional of the DGM with $L^2$-loss attempts to minimize the $L^2$-error of both the PDE and the terminal condition during training, ultimately learning the parameters that best approximate the solution of the PDE. In \cite{dgm}, the authors introduce the DGM, illustrate its ability to numerically solve high-dimensional nonlinear PDEs, and provide a convergence guarantee for second-order nonlinear parabolic PDEs.

The second popular deep learning-based method for solving high-dimensional PDEs, introduced in \cite{deep-bsde-orig} and expanded upon in \cite{nonlinear-feynman-kac, db-schemes}, exploits the connection between nonlinear parabolic PDEs and backward SDEs can be exploited via a so-called non-linear Feynman--Kac formula. In turn, the resulting backward SDE can be solved numerically by recursively using a sequence of neural networks to solve the SDE along a specified discretization of the time interval in question, starting with the terminal condition of the original PDE. Although this method, often referred to as deep BSDE, is likely applicable to our context, we defer further consideration of deep backward schemes to future work.

\subsection{Main Results}
\label{subsec:main-results}
We show that the convergence guarantee provided in \cite{dgm} extends to the (first-order) HJB equation \eqref{eq:mfcp-hjb}, relying on the theory of viscosity solutions for HJB equations to obtain the desired convergence. Importantly, our class of HJB equations does not fall into the class of second-order nonlinear parabolic PDEs considered in \cite{dgm}, so we provide a different proof technique. In order to leverage the theory of viscosity solutions in our proof technique, we modify the loss functional used in \cite{dgm} to approximate the uniform norm rather than the $L^2$-norm. Namely, we choose an $L^\infty$-loss. Indeed, viscosity solutions are often referred to as $L^\infty$-weak solutions, as they respect convergence with respect to the uniform norm \cite{crandall-evans-lions}. The work by \cite{dgm} establishes the uniform convergence of neural network approximators to the actual solution a family of partial differential equation. However, this result is obtained by introducing additional assumptions of uniform boundedness and equicontinuity on the neural networks, whch we bypass using the $L^\infty$-loss. With this reformulation of the DGM loss, we first prove an existence result by showing in Theorem~\ref{th:unif-error-conv} that a sequence of neural networks taking the loss functional of the DGM to zero exists. We then prove in Theorem~\ref{th:unif-value-conv} that for such a sequence of networks, the neural network approximators converge uniformly to the true value function for the MFCP on $[0, T] \times S_d$. 

To obtain this result, we utilize a standard argument from the theory of viscosity solutions, relying on an appropriate definition of viscosity solutions to Equation~\eqref{eq:mfcp-hjb} on the simplex and a corresponding comparison principle. By establishing that a sequence of neural network approximators are viscosity solutions to a sequence of perturbed HJB equations and applying the comparison principle to two suitable upper and lower limits of the neural network approximators, we can conclude that the sequence of neural network approximators taking the DGM loss functional to zero must converge uniformly to the classical solution to Equation~\eqref{eq:mfcp-hjb}. The main steps in our proof are outlined as follows:
\begin{itemize}
    \item First, in Theorem~\ref{th:unif-error-conv}, we show that there exists a neural network $\varphi$ that makes the DGM loss (defined in Section~\ref{sec:dgm}) arbitrarily small. 
    \item From this, in Corollary~\ref{corr:convergence-to-v}, we establish that there exists a sequence of neural network approximators $\{\varphi^n\}_{n \in \N}$ such that $\varphi^n \to V$ uniformly on $[0,T]\times S_d$ as $n \to \infty$, where $V$ is the unique classical solution to Equation~\eqref{eq:mfcp-hjb}.
    \item Then, we handle the technical details of working with a PDE on the $(d-1)$--dimensional simplex $S_d$. In particular, the interior of $S_d$ is empty, so we instead work on $\widehat{S}_d$, the projection of the simplex onto $\R^{d-1}$. These technical details are discussed in Proposition~\ref{prop:eps-n-convergence}, and we work on $\widehat{S}_d$ for the remainder of the proof in order to fully leverage the power of viscosity solutions. Throughout, we use the {\it hat} notation to denote mathematical objects that are defined on $\widehat{S}_d$ rather than on $S_d$. Specifically, $\widehat{V}$ is the version of $V$, defined on $\widehat{S}_d$.
    \item Finally, we arrive at our main result in Theorem~\ref{th:unif-value-conv}, which establishes that if $\{\widehat{\varphi}^n\}_{n \in \N}$ is a sequence of neural network approximators taking the DGM loss to zero (which we know exists from Corollary~\ref{corr:convergence-to-v}), then $\widehat{\varphi}^n \to \widehat{V}$ uniformly on $[0, T] \times \widehat{S}_d$. 
    \item The proof of Theorem~\ref{th:unif-value-conv} relies on a comparison principle for viscosity solutions, as presented in \cite[Theorem 3.3]{visc-users}. In particular, we construct a viscosity subsolution $\overline{V}$ and a viscosity supersolution $\underline{V}$ on $[0, T] \times \widehat{S}_d$ that satisfy $\underline{V} \leq \overline{V}$ by construction, and $\overline{V} \leq \underline{V}$ by the comparison principle. By the uniqueness result for viscosity solutions to the HJB equation \cite[Theorem 9]{mfcp}, this implies that $ \overline{V} = \underline{V} = \widehat{V}$, and the construction of $\overline{V}$ and $\underline{V}$ provides uniform convergence as a result. 
\end{itemize}

\subsection{Notation}
As noted above, we define $\lbrs{d} := \{1, \ldots, d\}$. The $(d-1)$--dimensional simplex $S_d$ is given by
\begin{align}
    \label{eq:simplex}
    S_d := \Big\{(m_1, \ldots, m_d) \in \R^{d} : m_i \geq 0, \quad \sum_{i=1}^d m_i = 1\Big\}.
\end{align}
On the simplex $S_d$, we define directional derivatives for a function $\varphi : S_d \to \R$ by
\begin{align*}
    D^i f(m) := (\partial_{m_j - m_i} f(m))_{j \in \lbrs{d}},
\end{align*}
where derivatives are in the direction of $e_j - e_i$ for the corresponding standard basis vectors of $\R^d$.

Unless otherwise specified, $\Omega$ will refer to a subset of $\R^n$ (measurable, but not necessarily open). As usual, $L^p(\Omega)$ consists of the measurable functions on $\Omega$ such that
\begin{align*}
    \|f\|_{p} := \left(\int_{\Omega} |f(x)|^p dx\right)^{1/p} < \infty,
\end{align*}
where the integral above is taken with respect to the Lebesgue measure on $\R^d$, restricted to $\Omega$. For $p = \infty$, we instead take $L^\infty(\Omega)$ to be the space of measurable functions on $\Omega$ whose essential supremum is bounded. We say that a function is differentiable on $\Omega$ if it is differentiable on the interior of $\Omega$ and extends to a differentiable function on an open neighborhood of $\R^n$ containing $\Omega$. As is standard, $\sC(\Omega)$ refers to the space of continuous functions on $\Omega$, equipped with the uniform norm:
\begin{align*}
    \|f\|_{\infty} := \sup_{x \in \Omega} |f(x)|.
\end{align*}
Below, when we refer to \textit{uniform convergence}, we always refer to convergence in the above norm.

We use the notation $\sC^{1,1}(\Omega)$ to denote (continuously-differentiable) functions that have Lipschitz-continuous first derivatives on $\Omega$. See \cite{mfcp} for a more thorough discussion of this class of functions as they relate to the HJB equation in Equation~\eqref{eq:mfcp-hjb}. Other spaces of functions that we refer to throughout the paper include the space of continuously-differentiable functions with $k$ continuously-differentiable derivatives, denoted by $\sC^k(\Omega)$. We define a norm on $\sC^k(\Omega)$ by
\begin{align*}
    \|f\|_{\sC^k(\Omega)} := \max_{|\alpha| \leq m} \sup_{x \in \Omega} |D^\alpha f(x)|.
\end{align*}

\subsection{Organization}
We structure the remainder of the paper as follows. Section~\ref{sec:hjb-mfcp-setup} introduces the relevant assumptions placed upon the MFCP that ensure sufficient regularity of the value function. Section~\ref{sec:dgm-conv} presents our modified version of the DGM algorithm for the HJB equation in Equation~\eqref{eq:mfcp-hjb}, the corresponding convergence results, and all associated proofs. Finally, Section~\ref{sec:numerical-tests} presents selected numerical results from a JAX implementation of the DGM.
\section{Assumptions and Known Results for the MFCP}
\label{sec:hjb-mfcp-setup}
Below, we outline the (already known) MFCP and the recent existence, uniqueness, and convergence results for the MFCP proven in \cite{mfcp}. We also present the relevant assumptions imposed upon the MFCP. This section utilizes the background on stochastic control and HJB equations provided in~\cite{stoch-control-text}. 



\subsection{Assumptions for the MFCP}
\label{subsec:assumptions}
We start by assuming the following, drawing upon some of the notations introduced throughout the preceding section.
\begin{assumption*}{Assumption A}
\label{a:mfcp-A}
The function $F : [0, T] \times [0, M]^{d^2} \times S_d \to \R$, given by
\begin{align*}
    F(t, a^1, \ldots, a^d, m) := \sum_{i \in \lbrs{d}} m_i f(t, i, a^i, m)
\end{align*}
is continuous and there exists a constant $C>0$, such that for any $t,s\in[0,T]$, $a\in[0,M]^d$, and $m,p\in S^d$,
\begin{align*}
    |F(t,a,m) - F(s,a,p)| &\leq  C(|t-s| + |m-p|).
\end{align*}
\end{assumption*}
\begin{assumption*}{Assumption B}
\label{a:mfcp-B}
For each $i \in \lbrs{d}$, the running cost $f$ is continuously differentiable in $a$, $\nabla_a f$ is Lipschitz continuous with respect to $m$, and there exists $\lambda > 0$ such that
\begin{align*}
    f(t, i, b,m) \geq f(t, i, a,m) + \langle \nabla_a f(t, i, a, m), b-a \rangle + \lambda |b-a|^2.
\end{align*}
\end{assumption*}
\begin{assumption*}{Assumption C}
\label{a:mfcp-C}
We have that:
\begin{enumerate}
\item[(C1)] Defining $G : S_d \to \R$ by
\begin{align*}
    G(m) := \sum_{i \in \lbrs{d}} m_i g^i(m),
\end{align*}
we have that
$F(\cdot, a, \cdot) \in \sC^{1,1}([0, T] \times S_d)$ and $G \in \sC^{1,1}(S_d)$.
\item[(C2)] The function
\begin{align*}
    [0,T] \times [0, \infty)^{d \times d} \times \intr(S_d) \ni (t,w,m) \mapsto \sum_{i \in \lbrs{d}}m_i f\Big(t, i, \Big(\frac{w_{i,j}}{m_i}\Big)_{j\neq i}, m\Big)
\end{align*}
is convex in $(w, m)$.
\item[(C3)] $G$ is convex in $m$.
\end{enumerate}
\end{assumption*}
Note that under Assumption~\hyperref[a:mfcp-C]{(C)}, $G$ is Lipschitz continuous: for any $m, p \in S_d$,
\begin{align*}
    |G(m) - G(p)| &\leq C|m-p|
\end{align*}
for some $C > 0$, which we can take to be the same constant as in Assumption~\hyperref[a:mfcp-A]{(A)}.


We provide an example of an MFCP in Section~\ref{sec:numerical-tests} for which the running and terminal costs satisfy all three of the above assumptions.
\subsection{Known Convergence, Existence, and Uniqueness results for the MFCP}
\label{subsec:prev-results}
Under assumptions similar to those in the preceding section, Cecchin \cite{mfcp} derived a series of useful results for the MFCP. To begin, Equation~\eqref{eq:mfcp-hjb} has a unique classical solution $V \in \sC^{1,1}([0,T] \times S_d)$.
\begin{proposition}[{{\cite[Theorem 2.9(3)]{mfcp}}}]
\label{prop:mfcp-hjb}
Under Assumptions~\hyperref[a:mfcp-A]{(A)}--~\hyperref[a:mfcp-C]{(C)} presented in Section~\ref{subsec:assumptions} above, the value function $V \in \sC^{1,1}([0, T] \times S_d)$ for the MFCP is the unique classical solution of the HJB equation \eqref{eq:mfcp-hjb}.
\end{proposition}

\begin{remark}
    If we only assume that Assumptions~\hyperref[a:mfcp-A]{(A)} and ~\hyperref[a:mfcp-B]{(B)} hold, then $V$ is both Lipschitz continuous in $(t, m)$ and the unique viscosity solution of Equation~\eqref{eq:mfcp-hjb} on $S_d$, and an optimal control for the MFCP exists \cite[Theorem 2.9(2)]{mfcp}. Under Assumption~\hyperref[a:mfcp-A]{(A)}, Cecchin also presents several convergence results connecting the MFCP to the $N$-agent optimization problem. In particular, we have, from \cite[Theorem 2.10]{mfcp}, the convergence of the value function $V^N$ for the $N$-agent optimization problem in \eqref{eq:n-agent} to the value function $V$ for the MFCP with a rate of $O(1/\sqrt{N})$.  Furthermore, \cite[Theorem 2.11]{mfcp} establishes the existence of an asymptotic optimal control for the $N$-agent optimization problem using the optimal control for the MFCP, also with an $O(1/\sqrt{N})$ rate of convergence.
\end{remark}

If in addition to Assumption~\hyperref[a:mfcp-A]{(A)}, we also include Assumptions~\hyperref[a:mfcp-B]{(B)}, then the empirical distribution corresponding to the optimal control for the MFCP, given in Equation~\eqref{eq:optimal-control-mfcp}, converges to the optimal process for the $N$-agent problem, with a convergence rate of $O(N^{-1/9})$ \cite[Theorem 2.13]{mfcp}.

\section{The DGM Algorithm and the Main Convergence Results}
\label{sec:dgm-conv}
In this section, we present our modified DGM algorithm to solve Equation \eqref{eq:mfcp-hjb} and the corresponding convergence proof. We begin by presenting our version of the DGM algorithm in Section~\ref{sec:dgm}. Next, in Section~\ref{sec:nn-approximators}, we describe the relevant results from universal approximation theory in addition to proving that the DGM can approximate the solution to Equation~\eqref{eq:mfcp-hjb} arbitrarily well in Theorem~\ref{th:unif-error-conv}. In Section~\ref{sec:nn-convergence}, we present our main convergence results in Theorem~\ref{th:unif-value-conv}. Finally, in Section~\ref{subsec:comp-original-dgm}, we describe the difficulties in establishing convergence for the algorithm with $L^2$-loss and explain with more details why the proof of \cite{dgm} is invalid for our equation \eqref{eq:mfcp-hjb}. We employ the following outline to obtain our existence and convergence results:
\begin{enumerate}
\item[(1)] By an appropriate version of the universal approximation theorem, we can arbitrarily approximate $V \in \sC^{1,1}([0,T] \times S_d)$ with neural network approximators; see Proposition~\ref{th:univ-approx}. Specifically, there exists a sequence of neural network approximators $\{\varphi(t, m; \theta^n)\}_{n \in \N}$ such that the DGM loss applied to the sequence converges to zero as $n \to \infty$; see Theorem~\ref{th:unif-error-conv}. 
\item[(2)] If the DGM loss goes to zero along this sequence of neural network approximators $\{\varphi(t, m; \theta^n)\}_{n \in \N}$, then $\varphi(\cdot,\cdot;\theta^n) \to V(\cdot,\cdot)$ uniformly on $[0, T] \times S_d$; see Theorem~\ref{th:unif-value-conv}.
\end{enumerate}
Throughout this section, we utilize the following operator, defined for $\phi \in \sC^{1,1}([0,T] \times S_d)$:
\begin{align}\label{eq:operator1}
    \sL[\phi](t,m) := -\partial_t \phi(t, m) + \sum_{i \in \lbrs{d}} m_i H^i(t, m, D^i \phi(t, m)).
\end{align}
Note that the first line in Equation \eqref{eq:mfcp-hjb} can be written as $\sL[V](t,m)=0$.
\subsection{DGM Algorithm}
\label{sec:dgm}
In this section, we present a modification of the DGM algorithm, first proposed by \cite{dgm}, in the context of the HJB equation for the MFCP. 
The DGM aims to efficiently approximate a solution to the above equation using a deep learning-based approach. Specifically, the method learns model parameters $\theta \in \R^P$, where $P$ depends on the dimension $d$ of the simplex $S_d$, the number of layers in the neural network in use, and the number of nodes in each layer of the neural network. The DGM learns the model parameters $\theta$ by minimizing an objective functional, referred to as the {\it DGM loss}, and given by 
\begin{align}
\label{eq:unif-loss}
\begin{split}
L(\theta) &:= \max_{(t, m) \in [0, T] \times S_d} |\sL[\varphi(\cdot, \cdot; \theta)](t,m)| + \max_{m \in S_d}|\varphi(T, m;\theta) - \sum_{i \in \lbrs{d}}m_i g^i(m)|. 
\end{split}
\end{align}
On occasions, we refer to this loss as the $L^\infty$-loss. 
The maxima over $[0, T] \times S_d$ and $S_d$ are approximated by sampling, as demonstrated in the following algorithm. As with the original DGM algorithm with $L^2$-loss \cite{dgm}, we utilize stochastic gradient descent (SGD) to find the parameter $\theta \in \R^P$ that minimizes the above loss. Note that in the following algorithm, the architecture of the neural network is fixed, and only the parameter $\theta$ is updated by SGD. As a proxy for the true loss functional, given a set of $M$ samples $(t^{(j)}, m^{(j)}, p^{(j)})_{j = 1, \ldots, M}:=\{(t^{(j)}, m^{(j)}, p^{(j)}):j = 1, \ldots, M\}  \subset [0,T] \times S_d\times S_d$, we define
\begin{align}
\label{eq:sampled-unif-loss}
\begin{split}
    G((t^{(j)}, m^{(j)}, p^{(j)})_{j = 1, \ldots, M}, \theta) &:= \max_{j = 1, \ldots, M}|\sL[\varphi(\cdot, \cdot; \theta)](t^{(j)},m^{(j)})| \\
    & \ + \max_{j = 1, \ldots, M} \Big|\varphi(T, p^{(j)};\theta) - \sum_{i \in \lbrs{d}}p_i^{(j)} g^i(p^{(j)})\Big|.
\end{split}
\end{align}
\begin{algorithm}[H]
\caption{Uniform DGM}\label{alg:unif-dgm}
\begin{algorithmic}
\State Initialize parameters $\theta^{(0)} \in \R^P$
\State Initialize tolerance $\delta \in (0,1)$
\State $n \gets 0$
\State Sample $(t^{(j)}, m^{(j)}, p^{(j)})_{j = 1, \ldots, M} \in [0,T] \times S_d$
\While{$G((t^{(j)}, m^{(j)}, p^{(j)})_{j = 1, \ldots, M}, \theta^{(n)}) \geq \delta$} 
\State Sample $(t^{(j)}, m^{(j)}, p^{(j)})_{j = 1, \ldots, M} \in [0,T] \times S_d$
\State $\theta^{(n+1)} \gets \theta^{(n)} - \alpha^{(n)} \nabla_\theta G((t^{(j)}, m^{(j)}, p^{(j)})_{j = 1, \ldots, M}, \theta^{(n)})$
\State $n \gets n + 1$
\EndWhile
\end{algorithmic}
\end{algorithm}
In practice, the performance of the above algorithm may vary depending on the sample size $M$ at each step. Additionally, the learning rate schedule $\alpha^{(n)}$ may determine the convergence rate of the algorithm as before. Using an optimizer such as AdaGrad or Adam may help speed up convergence. Finally, instead of using a tolerance $\delta \in (0,1)$ to determine the convergence of the algorithm, one may instead specify a fixed number of SGD iterations to carry out.

Although the DGM algorithm with $L^2$-loss as defined in \cite{dgm}, seems to work quite well in practice (as we demonstrate later in Section \ref{sec:numerical-tests}), the structure of the PDE in Equation~\eqref{eq:mfcp-hjb} prohibits us from using the same argument as \cite{dgm} to prove convergence of the DGM algorithm with $L^2$-loss. Instead, by slightly changing the loss function to the $L^\infty$-loss, given in \eqref{eq:unif-loss}, which we use to train the neural network approximation in the DGM algorithm, we can prove the convergence of our modified DGM algorithm to the unique value function of the MFCP. Specifically, by utilizing the above loss functional that approximates the uniform norm of the PDE and terminal condition rather than the squared error of the PDE and the terminal condition, we can leverage the theory of viscosity solutions for first-order HJB equations from \cite{visc-users} in our convergence proof. It is worth noting that the maximum is generally not a smooth function. In practice, the maximum is computed using a smooth maximum -- a function that approximates a maximum but is differentiable. For previous work involving stochastic gradient descent applied to a maximized loss function and insights into its robustness, refer to \cite{pmlr-v48-shalev-shwartzb16}. 
\begin{remark}
\normalfont
As noted above, we approximate that maxima over $[0, T] \times S_d$ and $S_d$ respectively by uniformly sampling $K$ points in each region. As shown in \cite[Lemma 1]{cohen2024deep}, we have that
\begin{align*}
    \E|L(\theta) - G((t^{(j)}, m^{(j)}, p^{(j)})_{j = 1, \ldots, M}, \theta)| = \mathcal{O}(|K|^{-1/(d-1)}),
\end{align*}
where $C_d > 0$ depends on the dimension $d$ and the parameters of Equation~\eqref{eq:mfcp-hjb}, and the expectation is taken with respect to the samples $(t^{(j)}, m^{(j)}, p^{(j)})_{j = 1, \ldots, M}$. Thus, by taking the sample size $K$ larger if necessary, we can accurately approximate \eqref{eq:unif-loss} via \eqref{eq:sampled-unif-loss}. Consequently, if the sampled loss converges to zero, then so does the true loss, up to an error term that depends on $d$, $K$, and the parameters of the HJB equation in question.
\end{remark}
\subsection{Neural Network Approximators}
\label{sec:nn-approximators} 
This section starts by utilizing the classical universal approximation result of Hornik \cite{univ-approx-thm} to establish that the uniform DGM loss can be made arbitrarily small by neural network approximators in Theorem~\ref{th:unif-error-conv}. Then, we lay the foundations for our main convergence result via Corollary~\ref{corr:convergence-to-v}, which provides a reformulation of the result in Theorem~\ref{th:unif-error-conv} that is better suited for the language of viscosity solutions. 

Although all universal approximation results in \cite{univ-approx-thm} hold for two-layer neural networks (i.e., a single hidden layer), we allow neural networks with multiple layers, including deep neural networks with modern architectures. In general, a network with $L$ layers, maximum width $n$, and a common activation function $\sigma$ takes the form
\begin{align}
\label{eq:multi-layer-nn}
    \varphi(t, m; \theta) := \sigma(W_{L} \ldots \sigma(W_1 m + \alpha t + c_1) \ldots + c_{L}),
\end{align}
where the activation function $\sigma$ is applied elementwise. Above, $W_i$ are weight matrices, $c_i$ are bias vectors, and $\alpha$ is a scalar weight. In turn, the parameters of each neural network are of the form $\theta = (W_1, \ldots, W_L, c_1, \ldots, c_L, \alpha) \in \R^P$ (upon flattening all weight matrices into vectors), where $P$ depends on the maximum width of the network, the depth $L$ of the network, and the dimension $d$ of the simplex $S_d$. In turn, we take $\mathfrak{C}^{(P)}_{d+1}(\sigma)$ to be the class of neural networks with parameters $\theta$ of dimension at most $P$, from which we define
\begin{align*}
    \mathfrak{C}_{d+1}(\sigma) := \bigcup_{P=1}^\infty \mathfrak{C}^{(P)}_{d+1}(\sigma).
\end{align*}
This is a slight departure from the notation of \cite{univ-approx-thm}, but all the relevant universal approximation results therein still hold in this more general context.
With the above notation in mind, we can apply the universal approximation theorem \cite[Theorem 3]{univ-approx-thm}, stated as follows:
\begin{proposition}
\label{th:univ-approx}
If $\sigma \in \sC^m(\R)$ is nonconstant and bounded, then $\mathfrak{C}_{d+1}$ is uniformly $m$-dense on compact sets in $\sC^m(\R^{d+1})$. In particular, for all $h \in \sC^m(\R^{d+1})$, all compact subsets $K \subset \R^{d+1}$, and any $\varepsilon > 0$, there exists $\psi = \psi(h, K, \varepsilon) \in \mathfrak{C}_{d+1}$ such that $\|h - \psi\|_{\sC^m(K)} < \varepsilon$.
\end{proposition}
In the implementation in Section~\ref{sec:numerical-tests}, we take $\sigma(y) = \tanh(y)$, a typical choice of activation function that is smooth, nonconstant, and bounded -- and therefore satisfies all the criteria of Proposition~\ref{th:univ-approx}. Note also that with this choice of $\sigma$, any element of $\mathfrak{C}_{d+1}(\sigma)$ is smooth (as a linear combination of smooth functions), ensuring that any $\varphi \in \mathfrak{C}_{d+1}(\sigma)$ has Lipschitz-continuous first derivative. See \cite{tanh-approx} for further justification of this choice of activation function in terms of the approximation guarantees that it brings.

With the above background in mind, we move towards approximating solutions to the HJB equation for the MFCP \eqref{eq:mfcp-hjb} using the DGM. Through assumptions~~\hyperref[a:mfcp-A]{(A)} --~\hyperref[a:mfcp-C]{(C)} in \cite{mfcp}, we have the following useful properties for proving convergence of the DGM to the solution of the HJB equation.
\begin{enumerate}
\item[(1)] Under Assumption~\hyperref[a:mfcp-A]{(A)}, the PDE-Hamiltonian in Equation~\eqref{eq:mfcp-hjb}, namely, $H(t, m, z)=\sum_{i \in \lbrs{d}} m_i H^i(t, m, z)$, is Lipschitz continuous in all three arguments. Specifically, we have that
\begin{align*}
    \Big|\sum_{i \in \lbrs{d}} m_i(H^i(t_1, m_1, z_1) - H^i(t_2, m_2, z_2))\Big| &\leq \kappa|(t_1 - t_2, m_1 - m_2, z_1 - z_2)| \\
    &\leq \kappa(|t_1 - t_2| + |m_1 - m_2| + |z_1 - z_2|)
\end{align*}
for some constant $\kappa > 0$, where $|\cdot|$ denotes the Euclidean norm. In fact, we only need the following, which also follows from the assumptions presented in \cite{mfcp}: for each $i \in \lbrs{d}$ and fixed $(t,m) \in [0,T] \times S_d$, the map $p \mapsto H^i(t, m, p)$ is Lipschitz continuous with a common Lipschitz constant $C > 0$.
\begin{proof}
By the definition of the Hamiltonian provided in \eqref{eq:hamiltonian}, we have that
\begin{align*}
    H^i(t, m, p) - H^i(t, m, p') &\leq \sup_{a_k \in [0, M], \;k\in\lbrs{d}\setminus\{i\}} \Big(-\sum_{k\in\lbrs{d}\setminus\{i\}}a_k p_k - f(t,i,a,m)\Big) \\
    & \ - \sup_{a_k \in [0, M], \;k\in\lbrs{d}\setminus\{i\}} \Big(-\sum_{k\in\lbrs{d}\setminus\{i\}}a_k p_k' - f(t,i,a,m)\Big) \\
    &= \sup_{a_k \in [0, M], \;k\in\lbrs{d}} \Big|\sum_{k\in\lbrs{d}\setminus\{i\}}a_k (p_k - p_k')\Big| \\
    &\leq M \Big|\sum_{k\in\lbrs{d}}(p_k - p_k')\Big| \\
    &\leq M |p - p'|,
\end{align*}
where $M$ is a bound on the transition rates as introduced in Section~\ref{subsec:mfcp}.  
Switching the roles of $p$ and $p'$ above completes the proof.
\end{proof}
\item[(2)] By Proposition \ref{prop:mfcp-hjb}, Equation~\eqref{eq:mfcp-hjb} admits a unique classical solution $V \in \sC^{1,1}([0,T] \times S_d)$.
\end{enumerate}

We show below that by utilizing the uniform error, given in Equation~\eqref{eq:unif-loss} 
we can obtain the desired convergence result. With the ultimate goal of showing that a neural network can approximate the value function $V(t, m)$ on $[0, T] \times S_d$ arbitrarily well in the uniform norm by taking
the number of neurons in the network, sufficiently large, we first show the following {\it existence result}. This result in fact holds for both the our DGM loss functional and the DGM with $L^2$-loss; see Appendix~\ref{app:a}. In particular, the following theorem establishes the existence of a sequence of neural networks that makes the DGM loss arbitrarily small.
\begin{theorem}
\label{th:unif-error-conv}
Let $\sigma \in \sC^1(\R)$ be bounded and nonconstant. For every $\varepsilon > 0$, there exists a constant $\Tilde{K}(d, T, C) > 0$, where $d$ is the dimension of the simplex $S_d$, $T$ is the finite time horizon of the MFCP, and $C$ is the Lipschitz constant of the PDE-Hamiltonian in Equation~\eqref{eq:mfcp-hjb}, such that for some $\varphi=\varphi(\cdot,\cdot;\theta) \in \mathfrak{C}_{d+1}(\sigma)$, the DGM loss functional in Equation~\eqref{eq:unif-loss} satisfies $L(\theta) \leq \Tilde{K} \varepsilon$.
\end{theorem}
\begin{proof}
Note that $\Omega_T := [0, T] \times S_d$ is a compact set in $\R^{d+1}$. Thus, by Proposition~\ref{th:univ-approx} above, we know that for the unique solution $V \in \sC^{1,1}(\Omega_T)$ to Equation~\eqref{eq:mfcp-hjb} and any $\varepsilon > 0$, there exists $\varphi \in \mathfrak{C}_{d+1}(\sigma)$ such that
\begin{align}
    \label{eq:univ-approx-est}
    \begin{split}
    \varepsilon > \sup_{(t,m) \in \Omega_T} |V(t,m) - \varphi(t, m; \theta)| &+ \sup_{(t,m) \in \Omega_T} |\partial_t V(t,m) - \partial_t \varphi(t, m; \theta)| \\
    &+ \sup_{(t, m) \in \Omega_T} |\nabla_m V(t,m) - \nabla_m \varphi(t,m; \theta)|  
    \end{split}
\end{align}
For such $\varphi \in \mathfrak{C}_{d+1}(\sigma)$, the Lipschitz continuity of the PDE-Hamiltonian in Equation~\eqref{eq:mfcp-hjb} yields
\begin{align*}
    \Big|\sum_{i \in \lbrs{d}} m_i H^i(t, m, D^i \varphi(t, m; \theta)) - \sum_{i \in \lbrs{d}} m_i H^i(t, m, D^i V(t, m))\Big| \leq C \sum_{i \in \lbrs{d}}|D^i \varphi(t, m; \theta) - D^iV(t,m)|,
\end{align*}
where $C$ is the Lipschitz constant of the PDE-Hamiltonian. Above, we use the fact that $|m_i| \leq 1$ for any $m \in S_d$. Next, for each $i \in \lbrs{d}$ and any $(t, m) \in \Omega_T$, we observe that, by denoting the standard basis of $\R^d$ by $\{e_i\}_{i \in \lbrs{d}}$, we have that
\begin{align}
\label{eq:gradient-error-estimate}
    |D^i \varphi(t, m; \theta) - D^iV(t,m)|^2 &= \frac{1}{2}\sum_{j=1}^d |(\nabla_m \varphi(t,m; \theta)- \nabla_m V(t,m)) \cdot (e_j - e_i)|^2 \notag \\
    &\leq \frac{1}{2}\sum_{j=1}^d |\nabla_m V(t,m) - \nabla _m\varphi(t,m; \theta)|^2 |e_j - e_i|^2 \notag \\
    &\leq d |\nabla_m V(t,m) - \nabla_m \varphi(t,m; \theta)|^2,
\end{align}
applying the Cauchy--Schwarz inequality in the second line above. Putting both bounds together, we obtain
\begin{align}
\label{eq:unif-error-estimate}
    \Big|\sum_{i \in \lbrs{d}} m_i H^i(t, m, D^i \varphi(t, m; \theta)) - \sum_{i \in \lbrs{d}} m_i H^i(t, m, D^i V(t, m))\Big| \leq C d^{3/2} |\nabla_m V(t,m) - \nabla_m \varphi(t,m; \theta)|.
\end{align}
Now, let $\Tilde{K} = C d^{3/2} > 0$. We conclude by noting that the value function $V$ satisfies $\sL[V](t,m) = 0$ for all $(t,m) \in \Omega_T$ (in addition to the terminal condition of Equation~\eqref{eq:mfcp-hjb}), which allows us to conclude that for $\varphi=\varphi(\cdot,\cdot;\theta)$,
\begin{align*}
    L(\theta) &= \max_{(t,m) \in \Omega_T}\left|\sL[\varphi](t,m)\right| + \max_{m \in S_d}\Big|\varphi(T, m; \theta) - \sum_{i \in \lbrs{d}} m_i g^i(m)\Big| \\
    &= \max_{(t,m) \in \Omega_T}\left|\sL[\varphi](t,m)\ - \sL[V](t,m)\right|+ \max_{m \in S_d}\left|\varphi(T, m; \theta) - V(T, m)\right|\\
    &\leq \max_{(t,m) \in \Omega_T}\Big|m_i H^i(t, m, D^i \varphi(t, m; \theta)) - \sum_{i \in \lbrs{d}} m_i H^i(t, m, D^i V(t, m))\Big| \\
    &\quad + \max_{(t,m) \in \Omega_T}\left|\partial_t V(t,m) - \partial_t \varphi(t,m;\theta)\right|\\
    &\quad + \max_{m \in S_d}\left|\varphi(T, m; \theta) - V(T, m)\right|\\
    &\leq \Tilde{K} \varepsilon
\end{align*}
by applying the approximation result from \eqref{eq:univ-approx-est} and taking $\Tilde{K} > 0$ larger if necessary. Note that the constant $\tilde{K}$ now depends on $d, C,$ and $T$ as claimed.
\end{proof}

The following corollary, which also relates to the existence of an approximating sequence of neural networks, utilizes the same universal approximation theorem as Theorem~\ref{th:unif-error-conv}. In particular, we can obtain a sequence of neural networks that satisfies a corresponding sequence of PDEs with a measurable error term that uniformly converges to zero.
\begin{corollary}
\label{corr:convergence-to-v}
There exists a sequence of 
parameters $\{\theta^n\}_{n \in \N}$, such that $\varphi^n \to V$ uniformly as $n \to \infty$, where $V$ is the unique classical solution to Equation~\eqref{eq:mfcp-hjb},
and $\varphi^n(t,m):=\varphi(t,m;\theta^n)$.
Given such $\varphi^n$, we define $e^n : [0, T] \times S^d \to \R$ as follows:
\begin{align}
\label{eq:corr-hjb}
\begin{split}
    &e^n(t, m):=\sL[\varphi^n](t, m) ,\qquad t\in[0,T),\\
    &e^n(T, m):=\varphi^n(T, m) - \sum_{i \in \lbrs{d}} m_i g^i(m) .
\end{split}
\end{align}
Then, $\|e^n\|_{\infty} \to 0$ as $n \to \infty$, noting that each $e^n$ depends on the neural network parameter $\theta^n$.
\end{corollary}
\begin{proof}
From Proposition~\ref{th:univ-approx}, we know that there exists a sequence of neural networks $\{\varphi^n\}_{n \in \N} \subset \mathfrak{C}_{d+1}(\sigma) \subset \sC^{1,1}([0,T] \times S_d)$, parametrized by a set of parameters $\{\theta^n\}_{n \in \N}$, such that \eqref{eq:univ-approx-est} holds for $\varepsilon = n^{-1}$. Thus, it immediately follows that $\|\varphi^n - V\|_{\infty} \to 0$ as $n \to \infty$,  yielding a sequence $\varphi^n$ of neural networks parametrized by $\theta^n$ that converges uniformly to $V$, the classical solution to Equation~\eqref{eq:mfcp-hjb}.
By construction, $\varphi^n$ satisfies the PDE in \eqref{eq:corr-hjb}. Furthermore, because each $e^n: [0,T] \times S_d \to \R$ is continuous on both $[0, T) \times S_d$ and $\{T\} \times S_d$, it is evidently measurable.

Finally, to see that $\|e^n\|_\infty \to 0$ as $n \to \infty$, we reuse many of the estimates from the proof of Theorem~\ref{th:unif-error-conv}. In particular, we can write
\begin{align*}
    e^n(t, m) &= -\partial_t \varphi^n(t, m) + \sum_{i \in \lbrs{d}} m_i H^i(t, m, D^i \varphi^n(t,m)) \\
    &= \partial_t V(t,m) -\partial_t \varphi^n(t, m) + \sum_{i \in \lbrs{d}} m_i\left[H^i(t, m, D^i \varphi^n(t,m))- H^i(t, m, D^i V(t,m))\right],
\end{align*}
for $(t,m) \in [0,T) \times S_d$, using the fact that $V$ solves Equation~\eqref{eq:mfcp-hjb}. By the fact that $m \in S_d$ and the Lipschitz continuity of $H^i$ for each $i \in \lbrs{d}$, we again have that
\begin{align*}
    \Big|\sum_{i \in \lbrs{d}} m_i\left[H^i(t, m, D^i \varphi^n(t,m))- H^i(t, m, D^i V(t,m))\right]\Big| &\leq C \sum_{i \in \lbrs{d}} |D^i \varphi^n(t, m) - D^iV(t,m)| \\
    &\leq 2 dC |\nabla_m V(t,m) - \nabla_m \varphi^n(t,m)|.
\end{align*}
Thus, we have that for all $(t,m) \in [0,T) \times S_d$
\begin{align*}
    |e^n(t,m)| \leq 2dC|\nabla_m V(t,m) - \nabla_m \varphi^n(t,m)| + |\partial_t V(t,m) -\partial_t \varphi^n(t, m)| \leq \frac{2dC}{n} + \frac{1}{n},
\end{align*}
by the construction of $\varphi^n$ from \eqref{eq:univ-approx-est}. Similarly, at the terminal time $t = T$, observe that
\begin{align*}
    |e^n(T,m)| = \Big|\varphi^n(T, m) - \sum_{i \in \lbrs{d}} m_i g^i(m)\Big| = |\varphi^n(T,m) - V(T, m)| \leq \frac{1}{n}.
\end{align*}
Thus, it follows that $\|e^n\|_{\infty} \to 0$ as $n \to \infty$ as claimed.
\end{proof}
\begin{remark}
\normalfont
The above sequence of neural networks satisfies several properties based on its construction. Specifically, note that by the converse of the Arzelà--Ascoli theorem, $\{\varphi^n\}_{n \in \N}$ is uniformly bounded and equicontinuous on $[0, T] \times S_d$. In fact, the construction of $\varphi^n$ implies that it is uniformly bounded with respect to the standard norm on $\sC^{1}([0, T] \times S_d)$.
\end{remark}
\subsection{Convergence of Neural Network Approximators to Value Function}
\label{sec:nn-convergence}
We now discuss the convergence of a sequence of neural network approximators $\varphi^n$ to $V$, the unique classical solution of the HJB equation in Equation \eqref{eq:mfcp-hjb}. With the ultimate goal of establishing \textit{uniform} convergence of the neural network approximators $\varphi^n$ to the value function $V$, we rely on the theory of viscosity solutions to first-order nonlinear PDEs. This powerful theory, developed by Crandall, Evans, and Lions in the 1980s for the explicit purpose of approaching HJB equations (which often lack classical, differentiable solutions) \cite{visc-users}, will allow us to relate the neural network approximators $\varphi^n(t,m)$ to the value function $V$ via a sequence of first-order nonlinear PDEs. Then, using a version of the comparison principle for viscosity solutions, we obtain the desired convergence.


Below, it will be useful to instead consider the HJB equation for $\eta \in \widehat{S}_d$, where
\begin{align*}
    \widehat{S}_d := \Big\{(\eta_1, \ldots, \eta_{d-1}) \in \R^{d-1} : \eta_j \geq 0 \; \text{for all} \; j = 1,\ldots, d-1, \; \sum_{j=1}^{d-1} \eta_j \leq 1 \Big\}.
\end{align*}
Then, the simplex $S_d$ can be expressed as
\begin{align*}
    S_d = \Big\{(\eta, \eta^{-d}) : \eta \in \widehat{S}_d, \; \eta^{-d} = 1 - \sum_{j=1}^{d-1} \eta_j\Big\}.
\end{align*}
In turn, any function $v \in \sC^1(S_d)$ induces a function $\widehat{v} \in \sC^1(\widehat{S}_d)$, given by $  \widehat{v}(\eta)=v(\eta, \eta^{-d})$. Furthermore, the gradient $\nabla_\eta \widehat{v}(\eta) = (\partial_{\eta_j}\widehat{v}(\eta))_{j=1}^{d-1}$ satisfies $\partial_{\eta_j}\widehat{v}(\eta) = \partial_{m_j - m_d} v(m)$. Following \eqref{eq:operator1}, we define an operator for functions defined on  $[0,T]\times \widehat{S}_d$ by
\begin{align}
\label{eq:mfcp-operator}
    \widehat{\sL}[\widehat{\phi}](t,\eta) := -\partial_t \widehat{\phi}(t, \eta) + \sum_{i \in \lbrs{d-1}} \eta_i \widehat{H}^i(t, \eta, \nabla_\eta \widehat{\phi}(t,\eta)) + \eta^{-d} \widehat{H}^d(t, \eta, \nabla_\eta \widehat{\phi}(t,\eta)).
\end{align}
Now, the solution to Equation~\eqref{eq:mfcp-hjb} can be written $V(t, \eta, \eta^{-d}) = \widehat{V}(t, \eta)$, where $\widehat{V} \in \sC^{1,1}(\widehat{S}_d)$ is the unique solution to the modified HJB equation 
\begin{align}
\label{eq:chart-hjb}
\begin{split}
    &\widehat{\sL}[\widehat{V}](t,\eta)= 0,\\
    &\widehat{V}(T, \eta) = \sum_{i \in \lbrs{d-1}} \eta_i g^i(\eta, \eta^{-d}) + \eta^{-d}g^d(\eta, \eta^{-d}).
\end{split}
\end{align}
Above, the modified Hamiltonians take inputs in $[0, T] \times \widehat{S}_d \times \R^{d-1}$ and are defined by
\begin{align*}
    \widehat{H}^i(t, \eta, p) = H^i(t, \eta, \eta^{-d}, p_1 - p_i, \ldots, p_{d-1} - p_i, -p_{i}) \quad \text{ and } \quad \widehat{H}^d(t, \eta, p) = H^d(t, \eta, \eta^{-d}, p, 0),
\end{align*}
for $i\in\lbrs{d-1}$. 
As shown in \cite{mfcp}, Equation~\eqref{eq:chart-hjb} has a unique solution $\widehat{V} \in \sC^{1,1}([0,T] \times \widehat{S}_d)$. Additionally, $\widehat{S}_d$ is a compact subset of $\R^{d-1}$, allowing us to apply the universal approximation theorem exactly as above, now on $\widehat{S}_d$. By the definitions of the modified Hamiltonians (in terms of the original Hamiltonians), we can also apply the exact same argument as above to obtain a result equivalent to our Theorem~\ref{th:unif-error-conv} on $\widehat{S}_d$. In particular, observe that for all $i \in \lbrs{d-1}$ and $p,p' \in \widehat{S}_d$, we have that
\begin{align*}
    |\widehat{H}^i(t, \eta, p) - \widehat{H}^i(t, \eta, p')|^2 &= |H^i(t, \eta, \eta^{-d}, p_1 - p_i, \ldots, p_{d-1} - p_i, -p_{i}) \\
    &\qquad - H^i(t, \eta, \eta^{-d}, p_1' - p_i', \ldots, p_{d-1}' - p_i', -p_{i}')|^2 \\
    &\leq C^2 \Big(\sum_{j=1}^{d-1} ((p_j - p_i) - (p_j' - p_i'))^2 + (p_i - p_i')^2 \Big)\\
    &\leq C^2 \Big(2\sum_{j=1}^{d-1} (p_j - p_j')^2 + (2(d-1) + 1)(p_i - p_i')^2\Big) \\
    &\leq C^2(2(d-1) + 1) \Big(\sum_{j=1}^{d-1} (p_j - p_j')^2 + (p_i - p_i')^2\Big) \\
    &= D^2 |p - p'|^2,
\end{align*}
where $D^2 = 2C^2(2(d-1) + 1)$. This shows that $\widehat{H}^i$ is Lipschitz continuous in $p$ with Lipschitz constant $D \geq C > 0$ for $i \in \lbrs{d-1}$, and we similarly observe that
\begin{align*}
    |\widehat{H}^d(t, \eta, p) - \widehat{H}^d(t, \eta, p')| = |H^d(t, \eta, \eta^{-d}, p, 0) - H^d(t, \eta, \eta^{-d}, p', 0)| \leq C|p - p'| \leq D|p - p'|
\end{align*}
so that $\widehat{H}^i$ is Lipschitz continuous in $p$, with common Lipschitz constant $D > 0$, for all $i \in \lbrs{d}$. From this, the proof of a modified version of Theorem~\ref{th:unif-error-conv}, now on $[0, T] \times \widehat{S}_d$, can proceed exactly as before. Note that the original value function defined on the simplex can be recovered via $V(t, \eta, \eta^{-d}) = \widehat{V}(t, \eta)$ for $\eta \in \widehat{S}_d$. 

Now, we can reframe the problem in terms of Equation~\eqref{eq:chart-hjb}, which possesses a unique classical solution $\widehat{V} \in \sC^{1,1}([0, T] \times \widehat{S}_d)$. Recall that $\widehat{S}_d \subset \R^{d-1}$ is the preimage of the simplex in $\R^d$ under the chart introduced in Equation~\eqref{eq:chart-hjb}. Working with Equation~\eqref{eq:chart-hjb} rather than Equation~\eqref{eq:mfcp-hjb} allows us to cite results from the theory of viscosity solutions that require the domain of the relevant PDE to be open; note that $\intr(\widehat{S}_d)$ is an open subset of $\R^{d-1}$ whereas $S_d$ has empty interior in $\R^d$. Additionally, the following result demonstrates that the convergence result on $\widehat{S}_d$ translates to $S_d$ without any issues.
\begin{proposition}
\label{prop:eps-n-convergence}
Assume that $\widehat{V} \in \sC([0,T] \times \widehat{S}_d)$ and $\widehat{\phi}^n \in \sC([0,T] \times \widehat{S}_d)$ are such that $\|\widehat{V} - \widehat{\phi}^n\|_\infty \to 0$ as $n \to \infty$. Then, $\|V - \phi^n\|_\infty \to 0$ as $n \to \infty$, where $V, \phi^n \in \sC([0,T] \times S_d)$ are given by $V(t, \eta, \eta^{-d}) = \widehat{V}(t, \eta)$ and $\phi^n(t, \eta, \eta^{-d}) = \widehat{\phi}^n(t, \eta)$ for all $(t, \eta, \eta^{-d}) \in [0,T] \times S_d$ and all $n \in \N$.
\end{proposition}
We remark that the opposite direction is also true; given a sequence of functions on $[0, T] \times S_d$ that converge uniformly to $V \in \mathcal{C}([0, T] \times S_d)$, the corresponding functions on $[0, T] \times \widehat{S}_d$ converge to $\widehat{V} \in \mathcal{C}([0, T] \times \widehat{S}_d)$, given by $\widehat{V}(t, \eta) = V(t, \eta, \eta^{-d})$ in our notation above. However, we do not require the converse of Proposition~\ref{prop:eps-n-convergence}, and the proof is analogous.
\begin{proof}
This is a simple consequence of the definition of $\widehat{S}_d$. Indeed, if $\|\widehat{V} - \widehat{\phi}^n\|_\infty \to 0$ as $n \to \infty$. Then, for any $\varepsilon > 0$, we have that for all $n \in \N$ sufficiently large,
\begin{align*}
    \sup_{(t,\eta) \in [0,T] \times \widehat{S}_d} |\widehat{V}(t, \eta) - \widehat{\phi}^n(t, \eta)| < \varepsilon.
\end{align*}
Consequently, for all $n \in \N$ sufficiently large, we have that
\begin{align*}
    \sup_{(t, \eta, \eta^{-d}) \in [0,T] \times S_d} |V(t,\eta, \eta^{-d}) - \phi^n(t,\eta, \eta^{-d})| &= \sup_{(t,\eta) \in [0,T] \times \widehat{S}_d} |\widehat{V}(t,\eta) - \widehat{\phi}^n(t,\eta)| \\
    &< \varepsilon.
\end{align*}
This means precisely that $\|V - \phi^n\|_\infty \to 0$ as $n \to \infty$.
\end{proof}
As a consequence of the above proposition, it suffices to show the uniform convergence of a sequence of neural network approximators $\widehat{\varphi}^n$ to the unique classical solution $\widehat{V}$ of Equation~\eqref{eq:chart-hjb} on $\widehat{S}_d$, as we can then recover uniform convergence on the simplex. Note that from now on, we consider neural networks, denoted by $\widehat{\varphi}$, which are defined as in \eqref{eq:multi-layer-nn} but with the input $\eta\in\widehat{S}_d$ rather than $m\in S_d$. From the discussion preceding Equation~\eqref{eq:chart-hjb}, we also know that Theorem~\ref{th:unif-error-conv} holds on $\widehat{S}_d$, yielding the existence of a sequence of neural networks $\{\widehat{\varphi}^n(t,\eta):=\widehat{\varphi}(t,\eta; \theta^n)\}_{n \in \N}$ such that
$L(\theta^n) \to 0$ as $n \to \infty$. In turn, each network $\widehat{\varphi}^n(t,\eta)$ satisfies its own ``perturbed'' PDE, of the form
\begin{align}
\label{eq:perturbed-hjb}
    \widehat{\sL}[\widehat{\varphi}^n](t, \eta) = \widehat{e^n}(t, \eta), \qquad (t,\eta) \in [0,T] \times \widehat{S}_d,
\end{align}
with $\widehat{e}^n(t, \eta):=e^n(t,\eta,\eta^{-d})$. 
For notational simplicity, we take
\begin{align*}
    \widehat{G}(\eta) := \sum_{i \in \lbrs{d-1}} \eta_i g^i(\eta, \eta^{-d}) + \eta^{-d}g^d(\eta, \eta^{-d})
\end{align*}
in this section to denote the terminal condition of the HJB equation on $\widehat{S}_d$. Denoting $\widehat{\Omega}_T := [0, T] \times \widehat{S}_d$ as in the previous section, Theorem~\ref{th:unif-error-conv} above implies that Equation~\eqref{eq:perturbed-hjb} satisfies
\begin{align*}
    \max_{(t,\eta) \in [0,T] \times \widehat{S}_d}\left|\widehat{e^n}(t,\eta)\right|+ \max_{\eta \in \widehat{S}_d}\left|\widehat{\varphi}^n(T, \eta) - \widehat{G}(\eta)\right| \to 0\qquad\text{as}\qquad n\to\infty.
\end{align*}
With this context in mind, we state the main {\it convergence result} of this section.
\begin{theorem}
\label{th:unif-value-conv}
The family of neural network approximators $\{\widehat{\varphi}^n(t, \eta)\}_{n \in \N}$ satisfying Equation~\eqref{eq:perturbed-hjb} converges uniformly to $\widehat{V} \in \sC^{1,1}([0,T] \times \widehat{S}_d)$, the unique classical solution of Equation~\eqref{eq:chart-hjb}, in the sense that
\begin{align*}
    \sup_{(t,\eta) \in [0,T] \times \widehat{S}_d} |\widehat{\varphi}^n(t, \eta) - \widehat{V}(t, \eta)| \to 0\qquad \text{as}\qquad n\to\infty.
\end{align*}
\end{theorem}
To prove the above theorem, we argue via the comparison principle for viscosity solutions to \eqref{eq:mfcp-hjb} presented in \cite{mfcp}. To this end, we require a suitable definition of viscosity solutions of Equation~\eqref{eq:chart-hjb} on $\widehat{S}_d$.
\begin{definition}
\label{def:visc-solutions}
A function $\widehat{v} \in \sC((0, T) \times \intr(\widehat{S}_d))$ is: 
\begin{enumerate}
\item[(i)] a viscosity subsolution of Equation~\eqref{eq:chart-hjb} if for any $\widehat{\varphi} \in \sC^1((0, T) \times \intr(\widehat{S}_d))$, $\widehat{\mathcal{L}}[\widehat{\varphi}] \leq 0$
for every local maximum $(t_0, \eta_0) \in (0, T) \times \intr(\widehat{S}_d)$ of $\widehat{v} - \widehat{\varphi}$ on $(0, T) \times \intr(\widehat{S}_d)$.
\item[(ii)] a viscosity supersolution of Equation~\eqref{eq:chart-hjb} if for any $\widehat{\varphi} \in \sC^1((0, T) \times \intr(\widehat{S}_d))$, $\widehat{\mathcal{L}}[\widehat{\varphi}] \geq 0$
for every local minimum $(t_0, \eta_0) \in (0, T) \times \intr(\widehat{S}_d)$ of $\widehat{v} - \widehat{\varphi}$ on $(0, T) \times \intr(\widehat{S}_d)$.
\item[(iii)] a viscosity solution of Equation~\eqref{eq:chart-hjb} if $\widehat{v}$ is both a viscosity subsolution and viscosity supersolution.
\end{enumerate}
\end{definition}
\begin{remark}
\normalfont
When viscosity solutions are introduced in \cite{mfcp}, the author also allows for test functions on $[0, T) \times S_d$ (resp. $[0, T) \times \widehat{S}_d$), noting that $[0, T) \times S_d$ (resp. $[0, T) \times \widehat{S}_d$) is no longer an open subdomain of $\R^{d+1}$ (resp. $\R^d$). However, in order to utilize \cite[Theorem 3.3]{visc-users}, the standard comparison principle for viscosity solutions, we must consider viscosity solutions on open subdomains of $\R^{d}$. As noted in \cite{mfcp}, it is also not immediately clear that a classical solution to Equation~\eqref{eq:mfcp-hjb} is a viscosity solution if the latter is defined on a closed set.

We could alternatively cite the comparison principle from \cite[Theorem 3.4]{mfcp} that utilizes the definition of viscosity solutions on closed sets presented therein. However, in order to utilize the clearly presented stability properties of viscosity solutions under uniform limits presented in \cite{ crandall-evans-lions, visc-users}, we opt for the standard definition in Definition~\ref{def:visc-solutions}.
\end{remark}
In order to establish Theorem~\ref{th:unif-value-conv}, we proceed using a standard comparison principle argument for viscosity solutions that also leverages the fact that $\widehat{V} \in \sC^{1,1}([0, T] \times \widehat{S}_d)$ is the unique viscosity solution to Equation~\eqref{eq:chart-hjb} from \cite{mfcp}.
\begin{proof}[Proof of Theorem~\ref{th:unif-value-conv}]
For each $n \in \N$, we may define an operator
\begin{align*}
    \widehat{\sL}^n[\widehat{\phi}](t, \eta) := -\partial_t \widehat{\phi}(t, \eta) + \sum_{i \in \lbrs{d-1}} \eta_i \widehat{H}^i(t, \eta, \nabla_\eta \widehat{\phi}(t,\eta)) + \eta^{-d} \widehat{H}^d(t, \eta, \nabla_\eta \widehat{\phi}(t,\eta)) - \widehat{e^n}(t,\eta),
\end{align*}
corresponding to the sequence of PDEs described in \eqref{eq:perturbed-hjb}. Because $\widehat{\sL}^n[\widehat{\phi}]$ depends only on the derivatives of $\widehat{\phi}$ (and not on $\widehat{\phi}$ itself), we observe that $\widehat{\sL}^n$ is \textit{proper} in the sense of \cite{visc-users}. This fact also ensures that the technical conditions preceding the comparison principle \cite[Theorem 3.3]{visc-users} are satisfied.

Now, note that by the discussion following Proposition~\ref{prop:eps-n-convergence},
$$\max_{(t,\eta) \in [0,T] \times \widehat{S}_d}\left|\widehat{e^n}(t,\eta)\right| \to 0\qquad \text{as}\qquad n\to\infty,$$
meaning that $\widehat{e^n}$ converges uniformly to zero on $[0,T] \times \widehat{S}_d$. Now, for each $n \in \N$, define $T^n : [0,T] \times \widehat{S}_d \times \R \times \R^{d-1} \to \R$ by
\begin{align*}
    T^n(t, \eta, p_0, p) := -p_0 + \sum_{i \in \lbrs{d-1}} \eta_i \widehat{H}^i(t, \eta, p) + p^{-d} \widehat{H}^d (t,\eta,p) - \widehat{e^n}(t, \eta).
\end{align*}
We then have that $T^n$ converges uniformly on $[0,T] \times \widehat{S}_d \times \R \times \R^{d-1}$ to
\begin{align*}
    T(t,\eta,p_0, p) := -p_0 + \sum_{i \in \lbrs{d-1}} \eta_i \widehat{H}^i(t, \eta, p) + p^{-d} \widehat{H}^d (t,\eta,p).
\end{align*}
These definitions are motivated by the fact that Equation~\eqref{eq:mfcp-hjb} can be written succinctly as
\begin{align*}
    T(t,\eta, \partial_t \widehat{\varphi}, \nabla_\eta \widehat{\varphi}) = 0,
\end{align*}
while Equation~\eqref{eq:perturbed-hjb} is given by
\begin{align*}
    T^n(t,\eta, \partial_t \widehat{\varphi}, \nabla_\eta \widehat{\varphi}) = 0
\end{align*}
for each $n \in \N$. Now, following \cite[Remark 6.3]{visc-users}, we note that because $\widehat{\varphi}^n$ is a classical solution (and therefore a viscosity solution) to the equation $T^n(t,\eta, \partial_t \widehat{\varphi}, \nabla_\eta \widehat{\varphi}) = 0$ on $(0,T) \times \intr(\widehat{S}_d)$, then 
\begin{align*}
    \overline{V}(t, \eta) := \lim_{j \to \infty} \sup\{\widehat{\varphi}^n(s, \nu) : n \geq j, \; (t, \eta) \in [0,T) \times \widehat{S}_d, \; |(s, \nu) - (t, \eta)| \leq 1/j\}
\end{align*}
is a viscosity subsolution to the equation $T(t,\eta, \partial_t \widehat{\varphi}, \nabla_\eta \widehat{\varphi}) = 0$, as we have that
\begin{align*}
    T(t,\eta,p_0,p) = \liminf_{n \to \infty} T^n(t,\eta,p,p_0).
\end{align*}
On the other hand, we also observe that 
\begin{align*}
    \underline{V}(t, \eta) := \lim_{j \to \infty} \inf\{\widehat{\varphi}^n(s,\nu) : n \geq j, \; (t, \eta) \in [0,T) \times \widehat{S}_d, \; |(s, \nu) - (t,\eta)| \leq 1/j\}
\end{align*}
is a viscosity supersolution to the equation $T(t,\eta, \partial_t \widehat{\varphi}, \nabla_\eta \widehat{\varphi}) = 0$ by the same reasoning. By construction, observe that $\underline{V} \leq \overline{V}$ on $[0,T) \times \widehat{S}_d$. Note also that both $\underline{V}$ and $\overline{V}$ are well-defined on $\{0\} \times \partial \widehat{S}_d$ by their construction. However, by the comparison principle presented in \cite[Theorem 3.3]{visc-users}, the fact that $\underline{V}$ is a viscosity supersolution and $\overline{V}$ is a viscosity subsolution is sufficient to conclude that $\overline{V} \leq \underline{V}$ on $[0,T) \times \widehat{S}_d$, observing that the comparison principle still holds on the closure of the domain $(0, T) \times \intr(\widehat{S}_d)$.

In particular, $\overline{V} = \underline{V}$ is a viscosity solution. As shown in \cite[Theorem 9]{mfcp}, Equation~\eqref{eq:mfcp-hjb} has a unique viscosity solution $\widehat{V}$, showing that $\overline{V} = \underline{V} = \widehat{V}$. Now, \cite[Remark 6.4]{visc-users} implies that $\lim_{n \to \infty} \widehat{\varphi}^n(t, \eta) = \widehat{V}(t, \eta)$ uniformly on $[0, T) \times \widehat{S}_d$.



Finally, note that we also have that
\begin{align*}
   \max_{\eta \in \widehat{S}_d}\left|\widehat{\varphi}^n(T, \eta) - \widehat{V}(T, \eta)\right| \to 0,
\end{align*}
from the construction of the modified DGM loss, allowing us to conclude uniform convergence of $\widehat{\varphi}^n \to \widehat{V}$ on the entire region $[0, T] \times \widehat{S}_d$ as claimed.
\end{proof}

\subsection{Comparison to the DGM with $L^2$-loss}
\label{subsec:comp-original-dgm}
\normalfont
At this point, we can clarify the reasons for the modification to the DGM algorithm made in Section~\ref{sec:dgm}. Here, we refer to the Sobolev space $W^{k,p}(\Omega)$, given by the space of functions $f \in L^p(\Omega)$ such that for any multi-index $\alpha$ with $|\alpha| \leq k$, the derivative $D^\alpha f$ exists and belongs to $L^p(\Omega)$ itself; see \cite[Chapter 5]{evans-pde} for more details on this characterization of Sobolev spaces. We also define the space $L^p(0, T; W^{k, q}(\Omega))$ by the function $f \in L^p([0, T] \times \Omega)$ such that, for fixed $t \in [0, T]$, $f(t, \cdot) \in W^{k, q}(\Omega)$. The spaces $L^p(0, T; L^q(\Omega))$ are then defined analogously for $1 \leq p, q \leq \infty$. 

Sirignano and Spiliopoulos \cite{dgm} formulated the $L^2$-loss,\footnote{As with the DGM loss in Equation~\eqref{eq:unif-loss}, this loss is computed in practice by sampling points according to probability measures $\nu_1$ and $\nu_2$ on $[0, T] \times S_d$ and $S_d$ respectively to obtain an unbiased estimate of the $L^2$-error: in Equation~\eqref{eq:l2-error}} as
\begin{align}
\label{eq:l2-error}
    \Tilde{L}(\theta) = \|\sL[\varphi(\cdot, \cdot; \theta)](t,m)\|_{2, [0, T] \times S_d, \nu_1}^2 + \|\varphi(T, m; \theta) - \sum_{i \in \lbrs{d}} m_i g^i(m)\|_{2, S_d, \nu_2}^2,
\end{align}
because of the natural connection between the class of equations that they considered and convergence in $L^2$. A key step in the proof of their analog to Theorem~\ref{th:unif-value-conv} involves obtaining a uniform bound on $\{\varphi(t, m; \theta^n)\}_{n \in \N}$ in $L^\infty(0, T; L^2(\Omega)) \cap L^2(0, T; W^{1,2}(\Omega))$, where $\Omega$ is the open domain on which the PDE is considered. In turn, this arises from an energy bound on quasilinear parabolic equations such as the one presented in \cite{energy-est-orig}, or in more generality in \cite{parabolic-est}. However, such a bound only holds for equations of the form
\begin{align*}
\begin{cases}
    \partial_t u - \text{div}(a(t,x, u, \nabla u)) = H(t,x, \nabla u) & (t,x) \in (0,T) \times \Omega, \\
    u = 0 & (t,x) \in (0,T) \times \partial \Omega, \\
    u(0,x) = u_0(x) & x \in \Omega.
\end{cases}
\end{align*}
that satisfy the Leray--Lions conditions. Namely, there must exist $\alpha > 0$ such that
\begin{align*}
    \alpha|\xi|^p \leq a(t,x, p, \xi) \cdot \xi
\end{align*}
for all $\xi \in \R^d$ and some $1 < p < d$. Clearly, this fails in our case, where $a$ is identically zero, even though our HJB equation otherwise satisfies the structure conditions in \cite{parabolic-est}.

Translating the convergence argument in \cite[Theorem 7.3]{dgm} to our context is also complicated by the fact that the class of quasilinear parabolic PDEs for which they prove convergence of the DGM possesses a standard notion of weak solutions that, via the dominated convergence theorem, cooperates with convergence in $L^2$. In the case of HJB equations, however, viscosity solutions take the place of weak solutions and instead behave nicely with respect to uniform convergence. Thus, we require that $\widehat{e^n} \to 0$ uniformly on $[0, T] \times \widehat{S}_d$ in Equation~\eqref{eq:perturbed-hjb}, but the formulation of the DGM with $L^2$-loss only implies convergence of the error term in $L^2$.

Finally, \cite{dgm} only concludes uniform convergence of the neural network approximators to the true solution of the PDE after imposing additional assumptions of uniform boundedness and equicontinuity on the neural networks. Although this may be a reasonable assumption to include, we find that our modified DGM algorithm and the theory of viscosity solutions provide a more direct route to uniform convergence; see Appendix~\ref{app:b} for a more detailed discussion of the equicontinuity of neural network approximators and the potential issues with such an approach.

Here, we also note an interesting connection between our work and recent work in the field of physics-informed neural networks (PINNs), a recent framework for incorporating PDE constraints into the training of neural networks. Specifically \cite{l2-pinns} show that, for a class of second-order HJB equations arising from stochastic control problems, it is \textit{impossible} to obtain a convergence guarantee along the lines of Theorem~\ref{th:unif-value-conv} if $L^2$-loss is used. We emphasize that this result does not necessarily apply to the first-order HJB equation arising from the MFCP. Our work is further complicated by the fact that Equation~\eqref{eq:mfcp-hjb} has solutions on the simplex $[0, T] \times S_d$, whereas the class of stochastic control problems considered in \cite{l2-pinns} are solved over $[0, T] \times \R^n$. However, future work may investigate whether the stability results obtained in \cite{l2-pinns} generalize to the context of the MFCP. Currently, this connection only provides some intuition as to why the $L^\infty$-loss may be the appropriate choice for training neural networks to solve Equation~\eqref{eq:mfcp-hjb}.
\section{Numerical Results}\label{sec:numerical-tests}
In this section, we present numerical results for the DGM applied to a simple example case of the MFCP, as presented in \cite[Example 2]{mfcp}. In particular, we work with the following example.
\begin{example}
\normalfont
Consider the quadratic running cost
\begin{align*}
    f(t, i, \alpha, m) = \frac{1}{2} \sum_{j\in\lbrs{d}, j \neq i} c_{i,j} \alpha_{i,j}^2 + f_0^i(m),
\end{align*}
with $f_0^i(m):= m_i$ and $\{c_{i,j}\}_{i,j \in \lbrs{d}} \in \R^{d \times d}$ a cost matrix that encodes the cost of transitioning from state $i$ to state $j$ for $i \neq j$. Finally, we consider the linear terminal condition given by
\begin{align*}
    g^i(m) = m_i
\end{align*}
for $i \in \lbrs{d}$. With this choice of terminal cost, we obtain the terminal condition
\begin{align*}
    V(T, m) = G(m) = \sum_{i \in \lbrs{d}} m_i^2.
\end{align*}
In this simple example, the Hamiltonian is explicitly given 
\begin{align*}
    H^i(t,m,z) = \sum_{j \neq i} \left(-\mathfrak{a}^*(-z_j)z_j - \frac{1}{2}(\mathfrak{a}^*(-z_j))^2\right) - f_0^i(m),
\end{align*}
where
\begin{align*}
    \mathfrak{a}^*(s) = \begin{cases}
    0 & s \leq 0, \\
    s & 0 \leq s \leq M, \\
    M & s \geq M.
    \end{cases}
\end{align*}
Recall that $M > 0$ is some constant such that $A = [0, M]^d$, where $A$ is the action space for the MFCP. Under this construction, all of the convexity and Lipschitz continuity constraints in Assumptions~\hyperref[a:mfcp-A]{(A)} -- \hyperref[a:mfcp-C]{(C)} are satisfied.

In the context of this example setup, we compare the performance of the DGM with $L^2$-loss from \cite{dgm} and the modified algorithm with our $L^\infty$-loss. Additionally, we explore the dependence of the performance of Algorithm~\ref{alg:unif-dgm} on the number of samples drawn from $[0, T] \times S_d \times S_d$. Finally, we demonstrate that the modified DGM algorithm scales well with dimension, providing numerical tests up to dimension $d = 200$. All numerical experiments were run on a Tesla T4 GPU via Google Colaboratory, and all code can be found in the \href{https://github.com/jakehofgard/jax-dgm}{GitHub repository} for this paper. 

For the sake of comparison, we utilize the same architecture as in \cite{dgm}; a neural network with four hidden LSTM-like layers and one dense output layer, with $\tanh$ activation throughout. We performed hyperparameter tuning on the number of layers, the number of samples used in each epoch, the width of the network, and the learning rate schedule. Unless otherwise noted, all results were produced by training the neural network for 200 epochs, with 10 gradient steps in each epoch, and $K = 10000$ samples for each epoch. After hyperparameter tuning, we found that a cosine one-cycle learning rate schedule, with a peak learning rate of $\alpha = 0.0008$, achieved the best performance. All experiments were performed with the Adam optimizer with weight decay and gradient clipping, which we found improved performance.

\subsection{Results and Discussion}
The plots in Figure~\ref{fig:d=2-unifsurface} and Figure~\ref{fig:d=2-l2surface} below demonstrate the value functions approximated by DGM with $L^\infty$-loss and $L^2$-loss respectively in dimension $d=2$. Both our DGM algorithm (with the $L^\infty$-loss) and the DGM algorithm with $L^2$-loss perform similarly, accurately solving the two-dimensional HJB equation. Figures~\ref{fig:l2-loss} and~\ref{fig:unif-loss} contain the corresponding loss curves. We remark that the loss curves for the two methods are displayed on different scales because they are \textit{not} comparable metrics of the performance of the solver. Given the formulation of the $L^\infty$-loss in Equation~\ref{eq:unif-loss}, we expect the $L^\infty$-loss to converge to a higher value than the $L^2$-loss, even if both metrics approximately solve the HJB equation.
\begin{figure}[H]
\centering
\begin{minipage}{.5\textwidth}
  \centering
  \includegraphics[width=\linewidth]{uniform_surface.png}
  \captionof{figure}{\footnotesize Approximate value function, learned using $L^\infty$-loss, in dimension $d = 2$.}
  \label{fig:d=2-unifsurface}
\end{minipage}%
\begin{minipage}{.5\textwidth}
  \centering
  \includegraphics[width=\linewidth]{l2_surface.png}
  \captionof{figure}{\footnotesize Approximate value function, learned using $L^2$-loss, in dimension $d = 2$.}
  \label{fig:d=2-l2surface}
\end{minipage}
\end{figure}

As demonstrated in Figures~\ref{fig:l2-loss} and~\ref{fig:unif-loss}, both metrics are susceptible to local minima using the LSTM-like architecture outlined in \cite{dgm}. Given sufficient training time and appropriate hyperparameter tuning, including the number of samples at each step, the learning rate schedule and the optimizer in use, both algorithms can closely approximate the true terminal condition of the example problem as demonstrated in Figure~\ref{fig:term-condition}. We expect that a deeper neural network, coupled with increased training time, would further improve accuracy. However, we defer such studies to a future work, having demonstrated the validity of the algorithm with a relatively simple architecture here.

\begin{figure}[H]
\centering
\begin{minipage}{.5\textwidth}
  \centering
  \includegraphics[width=\linewidth]{l2_loss.pdf}
  \captionof{figure}{\footnotesize $L^2$ training loss in dimension $d = 2$. PDE loss refers to the first term in the $L^2$-loss from Equation~\eqref{eq:l2-error}, terminal loss refers to the second term in Equation~\eqref{eq:l2-error}, and the combined loss represents the entirety of Equation~\eqref{eq:l2-error}.}
  \label{fig:l2-loss}
\end{minipage}%
\begin{minipage}{.5\textwidth}
  \centering
  \includegraphics[width=\linewidth]{uniform_loss.pdf}
  \captionof{figure}{\footnotesize $L^2$ training loss in dimension $d = 2$. PDE loss refers to the first term in the $L^\infty$-loss from Equation~\eqref{eq:unif-loss}, terminal loss refers to the second term in Equation~\eqref{eq:unif-loss}, and the combined loss represents the entirety of Equation~\eqref{eq:unif-loss}.}
  \label{fig:unif-loss}
\end{minipage}
\end{figure}

In Figure~\ref{fig:dep-samples-unif}, we demonstrate the relationship between the sample size $K$ and the stability of the loss for the uniform DGM algorithm. Unless otherwise specified, all numerical tests are carried out with $K = 10000$, as the tradeoff between stability and runtime becomes worse as the number of samples exceeds $K = 10000$. As $K$ increases, we observe more stable decreases in loss, as expected.
\begin{figure}[h]
\centering
\includegraphics[width=0.5\linewidth]{terminal_comparison.pdf}
\captionof{figure}{\footnotesize Approximate and true terminal conditions for both DGM algorithms in dimension $d = 2$.}
\label{fig:term-condition}
\end{figure}

Finally, in Table \ref{tab:runtimes}, we demonstrate the scalability of the DGM algorithm, implemented in JAX and run with GPU acceleration. By utilizing JAX's built-in auto-differentiation, just-in-time compilation, and GPU acceleration, the DGM algorithm scales well to dimension $d = 200$, with roughly linear increases in runtime past dimension $d = 100$. By runtime, we refer to the time that it takes to train the DGM network with uniform loss for $200$ epochs, with $10$ gradient descent steps in each epoch, and $K = 10000$ for each epoch. As the dimension increases, we observe that the loss decreases, as sampling points near the boundary of the simplex occurs with much lower probability in higher dimensions. Indeed, Figure~\ref{fig:term-condition} demonstrates that neither the $L^\infty$-loss nor the $L^2$-loss is able to fully learn the terminal condition near the boundary of the simplex, but this has a smaller impact on the loss in higher dimensions, thus resulting in lower losses as the dimension increases. However, as in dimension $d=2$, we still expect that the neural network learns the solution to the HJB equation relatively well away from the boundary of the simplex.
\begin{figure}[h]
\centering
\includegraphics[width=0.5\linewidth]{loss_uniform_error_samples.pdf}
\captionof{figure}{\footnotesize Dependence of $L^\infty$-loss on number of samples $K$. As the number of samples increases, the training process becomes more stable. Note that, with fewer samples, a lower loss does \textit{not} reflect a more accurate solution, as the sampled points are likely do not reflect the validity of the approximation on the entire simplex.}
\label{fig:dep-samples-unif}
\end{figure}
\begin{table}[H]
\caption{
\small Uniform DGM training times and losses as dimension $d$ increases. PDE loss refers to the first term in the DGM loss from Equation~\eqref{eq:unif-loss}, while terminal loss refers to the second term in Equation~\eqref{eq:unif-loss}.}
\label{tab:runtimes}
\vskip 0.15in
\begin{center}
\begin{small}
\begin{sc}
\begin{tabular}{lcccr}
\toprule
Dimension $d$ & Total Training Time ($s$) & Combined Loss & PDE Loss & Terminal Loss \\
\midrule
2    & 47.5 & 1.2134& 0.5452 &  0.6682 \\
5 & 56.1 & 0.7417 & 0.2508 & 0.4909 \\
10 & 62.0 & 0.7200 & 0.1190 & 0.6010 \\
20 & 65.0 & 0.2772 & 0.0392 &  0.2380 \\
50 & 76.0 & 0.0588 & 0.0093 & 0.0494 \\
100 & 100.0 & 0.0217 & 0.0031 & 0.0186 \\ 
200 & 221.0 & 0.0070 & 0.0011 & 0.0059 \\
\bottomrule
\end{tabular}
\end{sc}
\end{small}
\end{center}
\vskip -0.1in
\end{table}
\end{example}
\section{Conclusion}
We present a novel method for solving high-dimensional HJB equations on the simplex arising from MFCPs, complete with a convergence proof our algorithm. Our algorithm, which utilizes an $L^\infty$-loss functional for training rather than a standard $L^2$-loss functional, allows us to leverage the theory of viscosity solutions to prove that if a section of neural network approximators takes the $L^\infty$-loss functional to zero, then the sequence must uniformly converge to the unique viscosity solution of the HJB equation. Notably, our approach does \textit{not} assume boundedness and equicontinuity of the sequence of neural network approximators as in \cite{dgm}, and we thus provide a more general convergence guarantee than that of Sirignano and Spiliopoulos. In future work, we aim to relax our regularity assumptions for solutions to the HJB equation and investigate the convergence and performance of deep BSDE, another popular deep learning-based method for solving high-dimensional PDE, when applied to the MFCP.

\vspace{10pt}
\noindent
{\bf Acknowledgements.} 
The authors would like to Zihan Guo for his assistance with numerical experiments contained in this paper.

\bibliographystyle{unsrt}
\bibliography{refs}
\normalfont

\appendix
\section{Convergence Results for the DGM Algorithm with $L^2$-loss}
\label{app:a}
As discussed in Section~\ref{subsec:comp-original-dgm}, the proof technique via a comparison principle presented in Theorem~\ref{th:unif-value-conv} no longer holds for the DGM with $L^2$-loss, given by Equation~\eqref{eq:l2-error}. However, we \textit{do} still have an analog of Theorem~\ref{th:unif-error-conv} for the loss functional defined in Equation~\eqref{eq:l2-error}. For convenience, we recall that the DGM algorithm with $L^2$-loss aims to minimize the $L^2$-error of the approximate solution to the PDE in question. Specifically, as presented in \cite{dgm}, DGM learns an approximator $\varphi(t, m; \theta)$, parametrized by $\theta$, by minimizing the $L^2$-loss of the HJB equation:
\begin{align*}
    \Tilde{L}(\theta) := \|\sL[\varphi(\cdot, \cdot; \theta)](t,m)\|_{2, [0, T] \times S_d, \nu_1}^2 + \|\varphi(T, m; \theta) - \sum_{i \in \lbrs{d}} m_i g^i(m)\|_{2, S_d, \nu_2}^2.
\end{align*}
Above, $\nu_1$ and $\nu_2$ are probability densities on $[0, T] \times S_d$ and $S_d$ respectively. We then have the following analog to Theorem~\ref{th:unif-error-conv}:
\begin{theorem}
\label{th:l2-error-conv}
Let $\sigma \in \sC^1(\R)$ be bounded and nonconstant. For every $\varepsilon > 0$, there exists a constant $K(d, T, C) > 0$, where $d$ is the dimension of the simplex $S_d$, $T$ is the finite time horizon of the MFCP, and $C$ is the Lipschitz constant of the PDE-Hamiltonian in Equation~\eqref{eq:mfcp-hjb}, such that for some $\varphi=\varphi(\cdot,\cdot;\theta) \in \mathfrak{C}_{d+1}(\sigma)$, the DGM loss functional in Equation~\eqref{eq:l2-error} satisfies $\Tilde{L}(\theta) \leq K \varepsilon$.
\end{theorem}
\begin{proof}
The proof of this theorem is similar to that of Theorem~\ref{th:unif-error-conv}, with a few slight modifications. In particular, the Liptschitz continuity of the PDE-Hamiltonian yields
\begin{align}
    \label{eq:lip-bound-3.2}
    &\int_{\Omega_T} \Big|\sum_{i \in \lbrs{d}} m_i H^i(t, m, D^i \varphi(t, m; \theta)) - \sum_{i \in \lbrs{d}} m_i H^i(t, m, D^i V(t, m))\Big|^2 d\nu_1(t,m) \notag \\
    &\quad =\int_{\Omega_T} \Big|\sum_{i \in \lbrs{d}} m_i (H^i(t, m, D^i \varphi(t, m; \theta)) - H^i(t, m, D^i V(t, m)))\Big|^2 d\nu_1(t,m) \notag \\
    &\quad \leq d \sum_{i \in \lbrs{d}} \int_{\Omega_T} \left|m_i (H^i(t, m, D^i \varphi(t, m; \theta)) - H^i(t, m, D^i V(t, m)))\right|^2 d\nu_1(t,m)\notag \\
    &\quad \leq dC^2 \sum_{i \in \lbrs{d}} \int_{\Omega_T} |D^i \varphi(t, m; \theta) - D^iV(t,m)|^2 d\nu_1(t,m).
\end{align}
Above, we apply the Cauchy--Schwarz inequality in the second-to-last line above and note that $|m_i| \leq 1$ for any $m \in S_d$. Reusing the inequality from \eqref{eq:gradient-error-estimate}, we can bound \eqref{eq:lip-bound-3.2} by
\begin{align*}
    dC^2 \sum_{i \in \lbrs{d}} \int_{\Omega_T} |D^i \varphi(t, m; \theta) - D^iV(t,m)|^2 d\nu_1(t,m) &\leq d^3 C^2 \int_{\Omega_T} |\nabla_m (V(t,m) - \varphi(t,m; \theta))|^2d\nu_1(t,m) \\
    &\leq K \varepsilon^2
\end{align*}
for some positive constant $K = K(d, T, C) > 0$ by the construction of $\varphi$. Finally, because the value function $V$ satisfies $\sL[V](t,m) = 0$ for all $(t,m) \in \Omega_T$, we may write
\begin{align*}
    \Tilde{L}(\theta) &= \|\sL[\varphi](t,m)\|_{2, \Omega_T, \nu_1}^2 + \|\varphi(T, m; \theta) - \sum_{i \in \lbrs{d}} m_i g^i(m)\|_{2, S_d, \nu_2}^2 \\
    &= \|\sL[\varphi](t,m)\ - \sL[V](t,m)\|_{2, \Omega_T, \nu_1}^2 + \|\varphi(T, m; \theta) - V(T, m)\|_{2, S_d, \nu_2}^2 \\
    &\leq 2\int_{\Omega_T} \Big|\sum_{i \in \lbrs{d}} m_i H^i(t, m, D^i \varphi(t, m; \theta)) - \sum_{i \in \lbrs{d}} m_i H^i(t, m, D^i V(t, m))\Big|^2 d\nu_1(t,m) \\
    &\quad + 2\int_{\Omega_T} \left|\partial_t V(t,m) - \partial_t \varphi(t,m;\theta)\right|^2 d\nu_1(t,m) \\
    &\quad + \int_{S_d} |\varphi(T, m; \theta) - V(T, m)|^2 d\nu_2(m) \\
    &\leq K \varepsilon^2
\end{align*}
by applying the Cauchy--Schwarz inequality yet again, taking $K$ larger if necessary, and noting that the estimate in \eqref{eq:univ-approx-est} provides bounds on the two remaining terms in the above expression.
\end{proof}
\begin{remark}
\normalfont In the case of the DGM algorithm with $L^2$-loss, the measures $\nu_1$ and $\nu_2$, regardless of the densities that they correspond to, are defined as probability measures on $[0, T] \times S_d$ and $S_d$ respectively. Thus, the above result is \textit{independent} of the choice of densities $\nu_1$ and $\nu_2$, as we simply use the bounds
\begin{align*}
    \int_{\Omega_T} \left|\partial_t V(t,m) - \partial_t \varphi(t,m;\theta)\right|^2 d\nu_1(t,m) \leq \varepsilon^2 \nu_1(\Omega_T) = \varepsilon^2
\end{align*}
and
\begin{align*}
    \int_{S_d} |\varphi(T, m; \theta) - V(T, m)|^2 d\nu_2(m) \leq \varepsilon^2 \nu_2(S_d) = \varepsilon^2
\end{align*}
respectively.
\end{remark}
\section{Equicontinuity of Neural Network Approximators}
\label{app:b}
Recall that in \cite[Theorem 7.3]{dgm}, the authors require additional assumptions of both equicontinuity and uniform boundedness of neural network approximators that we bypass via the theory of viscosity solutions. The primary condition that allows for equicontinuity is the boundedness of the weights in the hidden layer(s) of the neural network used to approximate some continuous function. As shown in \cite{bdd-weights}, a general universal approximation theorem for continuous, bounded functions such as Proposition~\ref{th:univ-approx} above still holds for neural networks with bounded weights. In particular, it is likely possible to approximate $V \in \sC^{1,1}([0, T] \times S_d)$ (in $\sC^{1}$) via a sequence of neural networks with bounded weights. In turn, we may apply the equicontinuity results established in \cite{equi-nn}, obtaining an \textit{equicontinuous} sequence of neural network approximators to the value function $V$. Finally, the argument of \cite{dgm} provides for uniform boundedness, allowing us to apply the Arzelà--Ascoli theorem to obtain uniform convergence. This modification is stated more precisely below.
\begin{theorem}
\label{th:nn-equicont}
Take $\mathfrak{C}_{d+1}(\sigma)$ as defined above and consider any function $f \in \sC^{m}(K)$ for a compact set $K \subset \R^{d+1}$. Let $M > 0$ be such that $\sup_{x \in K} |f(x)| \leq M$. Now, let $\mathfrak{C}'_{d+1}(\sigma)$ denote the subset of networks in $\mathfrak{C}_{d+1}(\sigma)$ with weights $\theta = (\beta_1, \ldots, \beta_n, \alpha_{1,1}, \ldots, \alpha_{d+1,n}, c_1, \ldots, c_n) \in \R^{2n + n(d+1)}$ satisfying $|\alpha_{j, i}| \leq M$, $|c_i| \leq M$, and $|\beta_i| \leq M$ for all $i = 1,\ldots, n$ and $j = 1, \ldots, d+1$. Then, for any $m \in \N$, there exists a sequence of neural networks $\{\varphi^k(t,m):=\varphi(t,m;\theta^k)\}_{k \in \N}$ with an increasing number of hidden units such that $\|f - \varphi^k\|_{\sC^m(K)} < 1/k$ for all $k \in \N$. Furthermore, the sequence $\{\varphi^k\}_{k \in \N}$ is equicontinuous with respect to the inputs $(t,m) \in [0,T] \times S_d$.
\end{theorem}
\begin{proof}
As shown in \cite{bdd-weights}, any continuous, bounded functions can be uniformly approximated by neural networks with bounded weights. Assuming that we can obtain a universal approximation result for neural networks in $\mathfrak{C}'_{d+1}$, \cite[Proposition 8]{equi-nn} shows that the weights of any network $\varphi \in \mathfrak{C}'_{d+1}$ with $n$ hidden units such that $\|f - \varphi\|_{\infty} < 1$ must have weights satisfying
\begin{align*}
    \Big|\sum_{i=1}^n \beta_i\Big| \leq M + 1.
\end{align*}
In turn, the above condition allows us to apply \cite[Theorem 20]{equi-nn}, which states that the subset of $\mathfrak{C}_{d+1}'(\sigma)$ satisfying the above summability condition is equicontinuous with respect to the input space. Specifically, this implies that the sequence $\{\varphi^k\}_{k \in \N}$ is equicontinuous with respect to the inputs $(t,m) \in [0,T] \times S_d$ as claimed.
\end{proof}
In order to utilize Theorem~\ref{th:nn-equicont} in our context, however, we require a suitable version of the universal approximation theorem for neural networks with bounded weights. Despite a thorough literature review, it does not seem that the universal approximation result in Proposition~\ref{th:univ-approx} has an analogue for networks with uniformly bounded weights. In particular, there is no current result that allows for the approximation of an arbitrary function $h \in \sC^{m}(\R^{d+1})$ \textit{and} its derivatives on a compact set using networks with bounded weights. For the time being, the DGM algorithm with uniform loss and Theorem~\ref{th:unif-value-conv} provide a workaround for this issue, which we hope to resolve in future work.
\end{document}
\typeout{get arXiv to do 4 passes: Label(s) may have changed. Rerun}
% -------------------------
\begin{thebibliography}{10}

\bibitem{dgm}
Justin Sirignano and Konstantinos Spiliopoulos.
\newblock {DGM: A deep learning algorithm for solving partial differential equations}.
\newblock {\em Journal of Computational Physics}, 375:1339--1364, 2018.

\bibitem{drones7070476}
Min Yang, Guanjun Liu, Ziyuan Zhou, and Jiacun Wang.
\newblock Partially observable mean field multi-agent reinforcement learning based on graph attention network for uav swarms.
\newblock {\em Drones}, 7(7):476, 2023.

\bibitem{mot-pha2022CONTROL}
M\'{e}d\'{e}ric Motte and Huy\^{e}n Pham.
\newblock Mean-field {M}arkov decision processes with common noise and open-loop controls.
\newblock {\em Ann. Appl. Probab.}, 32(2):1421--1458, 2022.

\bibitem{7218597-2015}
Jian Li, Rajarshi Bhattacharyya, Suman Paul, Srinivas Shakkottai, and Vijay Subramanian.
\newblock Incentivizing sharing in realtime {D}2{D} streaming networks: A mean field game perspective.
\newblock In {\em 2015 IEEE Conference on Computer Communications (INFOCOM)}, pages 2119--2127, 2015.

\bibitem{crowd-motion-1}
Martin Burger, Marco Di~Francesco, Peter~A. Markowich, and Marie-Therese Wolfram.
\newblock On a mean field game optimal control approach modeling fast exit scenarios in human crowds.
\newblock In {\em 52nd IEEE Conference on Decision and Control}, pages 3128--3133, 2013.

\bibitem{crowd-motion-2}
Christian Dogbé.
\newblock Modeling crowd dynamics by the mean-field limit approach.
\newblock {\em Mathematical and Computer Modelling}, 52(9–10):1506–1520, November 2010.

\bibitem{duffy2010meanCONTROL}
Ken~R Duffy.
\newblock Mean field markov models of wireless local area networks.
\newblock {\em Markov Processes and Related Fields}, 16(2):295--328, 2010.

\bibitem{bor-sun2012CONTROL}
Vivek~S. Borkar and Rajesh Sundaresan.
\newblock Asymptotics of the invariant measure in mean field models with jumps.
\newblock {\em Stoch. Syst.}, 2(2):322--380, 2012.

\bibitem{aumann1964markets}
Robert~J. Aumann.
\newblock Markets with a continuum of traders.
\newblock {\em Econometrica}, 32:39--50, 1964.

\bibitem{schmeidler1973equilibrium}
David Schmeidler.
\newblock Equilibrium points of nonatomic games.
\newblock {\em J. Statist. Phys.}, 7:295--300, 1973.

\bibitem{lasry-lions-mfg}
Jean-Michel Lasry and Pierre-Louis Lions.
\newblock Mean field games.
\newblock {\em Japanese Journal of Mathematics}, 2(1):229--260, 2007.

\bibitem{huang-caines-malhame-mfg}
Minyi Huang, Roland~P. Malham{\'e}, and Peter~E. Caines.
\newblock {Large population stochastic dynamic games: closed-loop McKean-Vlasov systems and the Nash certainty equivalence principle}.
\newblock {\em Communications in Information \& Systems}, 6(3):221 -- 252, 2006.

\bibitem{Pham2017}
Huy\^en Pham and Xiaoli Wei.
\newblock Dynamic programming for optimal control of stochastic mckean--vlasov dynamics.
\newblock {\em SIAM Journal on Control and Optimization}, 55(2):1069–1101, January 2017.

\bibitem{Pham2018}
Huy\^en Pham and Xiaoli Wei.
\newblock Bellman equation and viscosity solutions for mean-field stochastic control problem.
\newblock {\em ESAIM: Control, Optimisation and Calculus of Variations}, 24(1):437–461, January 2018.

\bibitem{mfcp}
Alekos Cecchin.
\newblock Finite state {$N$-agent} and mean field control problems.
\newblock {\em ESAIM: COCV}, 27:31, 2021.

\bibitem{Bensoussan2013}
Alain Bensoussan, Jens Frehse, and Phillip Yam.
\newblock {\em Mean Field Games and Mean Field Type Control Theory}.
\newblock Springer New York, 2013.

\bibitem{Carmona2018_I}
René Carmona and Fran\c{c}ois Delarue.
\newblock {\em Probabilistic Theory of Mean Field Games with Applications I}.
\newblock Springer International Publishing, 2018.

\bibitem{Carmona2018_II}
René Carmona and Fran\c{c}ois Delarue.
\newblock {\em Probabilistic Theory of Mean Field Games with Applications II}.
\newblock Springer International Publishing, 2018.

\bibitem{ml-mfcp}
Lars Ruthotto, Stanley~J. Osher, Wuchen Li, Levon Nurbekyan, and Samy~Wu Fung.
\newblock A machine learning framework for solving high-dimensional mean field game and mean field control problems.
\newblock {\em Proceedings of the National Academy of Sciences}, 117(17):9183–9193, April 2020.

\bibitem{deep-learning-mfcp}
Jean-Pierre Fouque and Zhaoyu Zhang.
\newblock Deep learning methods for mean field control problems with delay.
\newblock {\em Frontiers in Applied Mathematics and Statistics}, 6:11, 2020.

\bibitem{mfc-mfcp-review}
René Carmona and Mathieu Laurière.
\newblock Deep learning for mean field games and mean field control with applications to finance, 2021.

\bibitem{scalable-rl-mfg}
Mathieu Laurière, Sarah Perrin, Sertan Girgin, Paul Muller, Ayush Jain, Theophile Cabannes, Georgios Piliouras, Julien Pérolat, Romuald Élie, Olivier Pietquin, and Matthieu Geist.
\newblock Scalable deep reinforcement learning algorithms for mean field games, 2022.

\bibitem{actor-critic-bsde}
Mo~Zhou, Jiequn Han, and Jianfeng Lu.
\newblock Actor-critic method for high dimensional static hamilton--jacobi--bellman partial differential equations based on neural networks.
\newblock {\em SIAM Journal on Scientific Computing}, 43(6):A4043–A4066, January 2021.

\bibitem{unique-optimal-control-mfg}
Diogo~A. Gomes, Joana Mohr, and Rafael~Rig{\~a}o Souza.
\newblock Continuous time finite state mean field games.
\newblock {\em Applied Mathematics \& Optimization}, 68(1):99--143, 2013.

\bibitem{kolokoltsov-conv}
Vassili N.~Kolokoltsov.
\newblock Nonlinear {Markov} games on a finite state space (mean-field and binary interactions).
\newblock {\em International Journal of Statistics and Probability}, 1(1):77--91, Apr 2012.

\bibitem{deep-bsde-orig}
Weinan E, Jiequn Han, and Arnulf Jentzen.
\newblock Deep learning-based numerical methods for high-dimensional parabolic partial differential equations and backward stochastic differential equations.
\newblock {\em Commun. Math. Stat.}, 5(4):349--380, 2017.

\bibitem{nonlinear-feynman-kac}
Huy\^{e}n Pham.
\newblock Feynman-{K}ac representation of fully nonlinear {PDE}s and applications.
\newblock {\em Acta Math. Vietnam.}, 40(2):255--269, 2015.

\bibitem{db-schemes}
C\^{o}me Hur\'{e}, Huy\^{e}n Pham, and Xavier Warin.
\newblock Deep backward schemes for high-dimensional nonlinear {PDE}s.
\newblock {\em Math. Comp.}, 89(324):1547--1579, 2020.

\bibitem{crandall-evans-lions}
M.~G. Crandall, L.~C. Evans, and P.~L. Lions.
\newblock {Some properties of viscosity solutions of Hamilton-Jacobi equations}.
\newblock {\em Transactions of the American Mathematical Society}, 282(2):487--502, 1984.

\bibitem{visc-users}
Michael~G. Crandall, Hitoshi Ishii, and Pierre-Louis Lions.
\newblock User's guide to viscosity solutions of second order partial differential equations.
\newblock {\em Bull. Amer. Math. Soc. (N.S.)}, 27(1):1--67, 1992.

\bibitem{stoch-control-text}
Huy\^{e}n Pham.
\newblock {\em Continuous-time stochastic control and optimization with financial applications}, volume~61 of {\em Stochastic Modelling and Applied Probability}.
\newblock Springer-Verlag, Berlin, 2009.

\bibitem{pmlr-v48-shalev-shwartzb16}
Shai Shalev-Shwartz and Yonatan Wexler.
\newblock Minimizing the maximal loss: How and why.
\newblock In Maria~Florina Balcan and Kilian~Q. Weinberger, editors, {\em Proceedings of The 33rd International Conference on Machine Learning}, volume~48 of {\em Proceedings of Machine Learning Research}, pages 793--801, New York, New York, USA, 20--22 Jun 2016. PMLR.

\bibitem{cohen2024deep}
Asaf Cohen, Mathieu Laurière, and Ethan Zell.
\newblock Deep backward and galerkin methods for the finite state master equation, 2024.

\bibitem{univ-approx-thm}
Kurt Hornik.
\newblock Approximation capabilities of multilayer feedforward networks.
\newblock {\em Neural Networks}, 4(2):251--257, 1991.

\bibitem{tanh-approx}
Tim {De Ryck}, Samuel Lanthaler, and Siddhartha Mishra.
\newblock On the approximation of functions by tanh neural networks.
\newblock {\em Neural Networks}, 143:732--750, 2021.

\bibitem{evans-pde}
Lawrence~C. Evans.
\newblock {\em Partial differential equations}, volume~19 of {\em Graduate Studies in Mathematics}.
\newblock American Mathematical Society, Providence, RI, second edition, 2010.

\bibitem{energy-est-orig}
Maria~Michaela Porzio.
\newblock Existence of solutions for some ``noncoercive'' parabolic equations.
\newblock {\em Discrete and Continuous Dynamical Systems}, 5(3):553 -- 568, 1999.

\bibitem{parabolic-est}
Martina Magliocca.
\newblock {Existence results for a Cauchy--Dirichlet} parabolic problem with a repulsive gradient term.
\newblock {\em Nonlinear Analysis}, 166:102--143, 2018.

\bibitem{l2-pinns}
Chuwei Wang, Shanda Li, Di~He, and Liwei Wang.
\newblock {Is $L^2$ Physics Informed Loss Always Suitable for Training Physics Informed Neural Network?}
\newblock In {\em Advances in Neural Information Processing Systems}, volume~35, pages 8278--8290. Curran Associates, Inc., 2022.

\bibitem{bdd-weights}
M.~Stinchcombe and H.~White.
\newblock Approximating and learning unknown mappings using multilayer feedforward networks with bounded weights.
\newblock In {\em 1990 IJCNN International Joint Conference on Neural Networks}, pages 7--16 vol.3, 1990.

\bibitem{equi-nn}
P.~Chandra and Y.~Singh.
\newblock Feedforward sigmoidal networks - equicontinuity and fault-tolerance properties.
\newblock {\em IEEE Transactions on Neural Networks}, 15(6):1350--1366, 2004.

\end{thebibliography}
